% Options for packages loaded elsewhere
\PassOptionsToPackage{unicode}{hyperref}
\PassOptionsToPackage{hyphens}{url}
%
\documentclass[
]{article}
\usepackage{amsmath,amssymb}
\usepackage{iftex}
\ifPDFTeX
  \usepackage[T1]{fontenc}
  \usepackage[utf8]{inputenc}
  \usepackage{textcomp} % provide euro and other symbols
\else % if luatex or xetex
  \usepackage{unicode-math} % this also loads fontspec
  \defaultfontfeatures{Scale=MatchLowercase}
  \defaultfontfeatures[\rmfamily]{Ligatures=TeX,Scale=1}
\fi
\usepackage{lmodern}
\ifPDFTeX\else
  % xetex/luatex font selection
\fi
% Use upquote if available, for straight quotes in verbatim environments
\IfFileExists{upquote.sty}{\usepackage{upquote}}{}
\IfFileExists{microtype.sty}{% use microtype if available
  \usepackage[]{microtype}
  \UseMicrotypeSet[protrusion]{basicmath} % disable protrusion for tt fonts
}{}
\makeatletter
\@ifundefined{KOMAClassName}{% if non-KOMA class
  \IfFileExists{parskip.sty}{%
    \usepackage{parskip}
  }{% else
    \setlength{\parindent}{0pt}
    \setlength{\parskip}{6pt plus 2pt minus 1pt}}
}{% if KOMA class
  \KOMAoptions{parskip=half}}
\makeatother
\usepackage{xcolor}
\usepackage[margin=1in]{geometry}
\usepackage{color}
\usepackage{fancyvrb}
\newcommand{\VerbBar}{|}
\newcommand{\VERB}{\Verb[commandchars=\\\{\}]}
\DefineVerbatimEnvironment{Highlighting}{Verbatim}{commandchars=\\\{\}}
% Add ',fontsize=\small' for more characters per line
\usepackage{framed}
\definecolor{shadecolor}{RGB}{248,248,248}
\newenvironment{Shaded}{\begin{snugshade}}{\end{snugshade}}
\newcommand{\AlertTok}[1]{\textcolor[rgb]{0.94,0.16,0.16}{#1}}
\newcommand{\AnnotationTok}[1]{\textcolor[rgb]{0.56,0.35,0.01}{\textbf{\textit{#1}}}}
\newcommand{\AttributeTok}[1]{\textcolor[rgb]{0.13,0.29,0.53}{#1}}
\newcommand{\BaseNTok}[1]{\textcolor[rgb]{0.00,0.00,0.81}{#1}}
\newcommand{\BuiltInTok}[1]{#1}
\newcommand{\CharTok}[1]{\textcolor[rgb]{0.31,0.60,0.02}{#1}}
\newcommand{\CommentTok}[1]{\textcolor[rgb]{0.56,0.35,0.01}{\textit{#1}}}
\newcommand{\CommentVarTok}[1]{\textcolor[rgb]{0.56,0.35,0.01}{\textbf{\textit{#1}}}}
\newcommand{\ConstantTok}[1]{\textcolor[rgb]{0.56,0.35,0.01}{#1}}
\newcommand{\ControlFlowTok}[1]{\textcolor[rgb]{0.13,0.29,0.53}{\textbf{#1}}}
\newcommand{\DataTypeTok}[1]{\textcolor[rgb]{0.13,0.29,0.53}{#1}}
\newcommand{\DecValTok}[1]{\textcolor[rgb]{0.00,0.00,0.81}{#1}}
\newcommand{\DocumentationTok}[1]{\textcolor[rgb]{0.56,0.35,0.01}{\textbf{\textit{#1}}}}
\newcommand{\ErrorTok}[1]{\textcolor[rgb]{0.64,0.00,0.00}{\textbf{#1}}}
\newcommand{\ExtensionTok}[1]{#1}
\newcommand{\FloatTok}[1]{\textcolor[rgb]{0.00,0.00,0.81}{#1}}
\newcommand{\FunctionTok}[1]{\textcolor[rgb]{0.13,0.29,0.53}{\textbf{#1}}}
\newcommand{\ImportTok}[1]{#1}
\newcommand{\InformationTok}[1]{\textcolor[rgb]{0.56,0.35,0.01}{\textbf{\textit{#1}}}}
\newcommand{\KeywordTok}[1]{\textcolor[rgb]{0.13,0.29,0.53}{\textbf{#1}}}
\newcommand{\NormalTok}[1]{#1}
\newcommand{\OperatorTok}[1]{\textcolor[rgb]{0.81,0.36,0.00}{\textbf{#1}}}
\newcommand{\OtherTok}[1]{\textcolor[rgb]{0.56,0.35,0.01}{#1}}
\newcommand{\PreprocessorTok}[1]{\textcolor[rgb]{0.56,0.35,0.01}{\textit{#1}}}
\newcommand{\RegionMarkerTok}[1]{#1}
\newcommand{\SpecialCharTok}[1]{\textcolor[rgb]{0.81,0.36,0.00}{\textbf{#1}}}
\newcommand{\SpecialStringTok}[1]{\textcolor[rgb]{0.31,0.60,0.02}{#1}}
\newcommand{\StringTok}[1]{\textcolor[rgb]{0.31,0.60,0.02}{#1}}
\newcommand{\VariableTok}[1]{\textcolor[rgb]{0.00,0.00,0.00}{#1}}
\newcommand{\VerbatimStringTok}[1]{\textcolor[rgb]{0.31,0.60,0.02}{#1}}
\newcommand{\WarningTok}[1]{\textcolor[rgb]{0.56,0.35,0.01}{\textbf{\textit{#1}}}}
\usepackage{graphicx}
\makeatletter
\def\maxwidth{\ifdim\Gin@nat@width>\linewidth\linewidth\else\Gin@nat@width\fi}
\def\maxheight{\ifdim\Gin@nat@height>\textheight\textheight\else\Gin@nat@height\fi}
\makeatother
% Scale images if necessary, so that they will not overflow the page
% margins by default, and it is still possible to overwrite the defaults
% using explicit options in \includegraphics[width, height, ...]{}
\setkeys{Gin}{width=\maxwidth,height=\maxheight,keepaspectratio}
% Set default figure placement to htbp
\makeatletter
\def\fps@figure{htbp}
\makeatother
\setlength{\emergencystretch}{3em} % prevent overfull lines
\providecommand{\tightlist}{%
  \setlength{\itemsep}{0pt}\setlength{\parskip}{0pt}}
\setcounter{secnumdepth}{-\maxdimen} % remove section numbering
\usepackage{booktabs}
\usepackage{longtable}
\usepackage{array}
\usepackage{multirow}
\usepackage{wrapfig}
\usepackage{float}
\usepackage{colortbl}
\usepackage{pdflscape}
\usepackage{tabu}
\usepackage{threeparttable}
\usepackage{threeparttablex}
\usepackage[normalem]{ulem}
\usepackage{makecell}
\usepackage{xcolor}
\ifLuaTeX
  \usepackage{selnolig}  % disable illegal ligatures
\fi
\usepackage{bookmark}
\IfFileExists{xurl.sty}{\usepackage{xurl}}{} % add URL line breaks if available
\urlstyle{same}
\hypersetup{
  pdftitle={UntitledHuman Supervision of Multi-Agent Reinforcement-Learning Teams: Selecting Teams across Environments with Varying Risk of Stag Illness},
  pdfauthor={Student name:},
  hidelinks,
  pdfcreator={LaTeX via pandoc}}

\title{UntitledHuman Supervision of Multi-Agent Reinforcement-Learning
Teams: Selecting Teams across Environments with Varying Risk of Stag
Illness}
\author{Student name:}
\date{2026-01-11}

\begin{document}
\maketitle

\section{Abstract}\label{abstract}

This paper has analyzed the selection of one out of three pre-trained
multi-agent reinforcement-learning teams (Stag, Hare, DisAdv) by human
supervisors in presented environments that differ in the likelihood of
the presence of high-reward stags who are sick (0, 30, 90). Each of the
thirty human supervisors was taken through nine ecology trials and gave
workload (NASA-TLX), situation awareness (SART), and
decision/behavioural measures (choice accuracy, response latency,
team-selection clicks, and performance score). The descriptive statistic
reveals that the mean performance grades are 691.2 ( Extreme Hare, E),
751.5 ( Ambiguous, O) and 795.6 ( Extreme Stag, R). One-way ANOVA showed
that there was a significant difference between condition and
performance as F(2,27) = 4.49, p = 0.0207, and post-hoc tests showed
that significantly better performance was observed in R than in E (mean
difference = 104.4, p = 0.0159). There was also a difference in
situation awareness (SART), F(2,27) = 5.24, p = 0.012, O \textgreater{}
E (diff = 4.9, p = 0.0177) and R \textgreater{} E (diff = 4.4, p =
0.0352). There was no significant difference between condition in
workload (NASA-TLX total), F(2,27) =1.98, p=0.158; Bayesian ANOVA
provided BF=0.733 (evidence slightly favouring the null). Correlational
results depict performance to have positive relationship with choice
accuracy ( r = 0.64, p = 0.001) and negative relationship with choice
accuracy response time ( r = -0.446, p = 0.014). Findings suggest
supervisors would make the best decisions and make them quicker in
unambiguous Stag-favouring conditions; ambiguous conditions would cause
more reported SA, but no higher choice of the adaptive DisAdv team. It
has implications to human-agent teaming, transparency and training.

\section{1. Introduction}\label{introduction}

In contemporary socio-technical systems, human supervisors should
choose, control, and supervise autonomous agents in which the learned
behaviour of the agents can vary significantly depending on the context
of training (Vaccaro, Almaatouq \& Malone, 2021; Schmutz et al., 2023).
Multi-agent reinforcement learning (MARL) is a model where agents
trained in an environment where they are directed to pursue either
rewards (e.g.~catching stags) or avoidance (e.g.~avoiding stags when
they are often sick) devises specialised strategies which are only
suitable in particular environmental situations. It is possible that a
hybrid or adaptive team (in this case called DisAdv) which is able to
sense the changing conditions and exchange information among agents can
be useful in ambiguous environments (Onnasch et al., 2014; Schmid and
Stanton, 2019). Nevertheless, successful human supervision is not based
solely on agent capacity but also on the situation awareness (SA) of the
human, perceived workload and capacity to interpret (and trust) agent
behaviour (Endsley, 1995; Hart and Staveland, 1988). This paper
discusses the skills of human supervisors to choose the best team (Stag,
Hare, DisAdv) in the environment with different stag-illness likelihood
(0\%, 30, 90). The experimental variables were: task performance (mark),
the correctness of the choice (proprobability of selecting team to be
the optimal team as per the environment), response latency (time to
select a team after the environment change), team-selection activity,
perceived workload (NASA-TLX subscales and total), and situation
awareness (SART). The design enabled us to test the hypothesis: (a)
supervisors will perform better in clear (extreme) vs ambiguous
environments; (b) supervisors will choose adaptive teams (DisAdv) in an
ambiguous environment; and (c) supervisors will have predictable
relationships between SA, workload and decision performance.

\subsection{1.1 Research Questions and
Hypotheses}\label{research-questions-and-hypotheses}

Based on the literature above, we were able to develop the following
testable hypotheses: I. \textbf{H1:} The accuracy and performance of
choices will be better in extreme environments (0\% and 90\%
stag-illness) than in the ambiguous environment (30\%).\\
II. \textbf{H2:} The adaptive DisAdv team will be more frequently chosen
by supervisors in the ambiguous environment compared to extreme
environments.\\
III. \textbf{H3:} In ambiguous environments, the selection latency (time
to click team selection) will be extended as compared to extreme
environments.\\
IV. \textbf{H4:} The accuracy of choice and short response times will
have a positive and negative relationship respectively with higher
situation awareness (SART scores).\\
V. \textbf{H5:} The self-reported workload (NASA-TLX total) will be
higher in the ambiguous setting and will have a negative relationship
with choice accuracy.\\
These hypotheses are based on the cognitive-systems studies that
indicate that clarity of environmental cues decreases the demands on
processing and enhances monitoring (Onnasch et al., 2014) and ambiguity
intensifies the necessity of adaptive strategies and imposes more
demands on SA (Endsley, 1995; Schmid and Stanton, 2019).

\section{2. Methods}\label{methods}

\subsection{2.1 Participants}\label{participants}

The study adopted within-subjects design, and every participant was
subjected to all three types of environment, and there were three
replicas of each. The defined environment conditions were 0\% of Stag
healthy and 30 percent of ambiguous and 90 percent of Stag ill.~The
participants had to choose one of three pre-trained AI teams, in each of
the three environments: the Stag team that employed a reward-seeking
policy; the Hare team that was risk-averse; and the DisAdv team that was
an adaptive team where one agent could detect and share stag illness
information. Before the task commenced the participants were given
instruction as to the workflow and the rules that applied to the
environments so that they were aware of how the environment changes are
reflected in the best team decisions.

\subsection{2.2 Experimental Design and
Materials}\label{experimental-design-and-materials}

The experiment operated using within subject design, that is, each
subject was subjected to all the three types of environments and three
replicas of each type were used (Lambdin and Shaffer, 2009). The
conditions of the environment were characterized as 0 percent of Stag
healthy, 30 per cent of ambiguous, and 90 percent of Stag ill.~In each
of the three environments, participants were needed to choose one of
three pre-trained AI teams, including the Stag team, which pursued a
reward-seeking strategy, the Hare team, which avoided risks, and the
DisAdv team, which was an adaptive team where one agent could learn
about stag illness and exchange information on it. The participants were
pre-informed of the workflow and the principles which governed the
environments, so that before starting with the task, they knew the
effect that environment manipulation had on the optimal team decisions
(Collins et al., 2022). A number of outcome measures were observed
during the task. Performance score was the numeric award that each
participant received in the exposure to the environment, and was put
down in FullData sheet and FullData2 sheet (Koopmans et al., 2014). The
measure of the accuracy of the choice consisted in the percentage of the
team selections that corresponded to the best team in the environment
0\% environments were best handled by the Stag team, 30\% by the DisAdv
team and 90\% by the Hare team. Latency was measured by Time to respond
to the environment which was the number of seconds in which the change
in environment occurred to the time the participant selected a team
(Emmanuelle Apartis, 2013). The team selection clicks reported the
number of times the participants interacted with the selection button
which could be checking or switching teams. The NASA-TLX was used to
measure subjective workload, and it has six subscales (mental demand,
physical demand, temporal demand, performance, effort, and frustration)
and a total score that was calculated as the mean of the subscales. The
SART was used to test the situation awareness which produced a total
score per participant. All data were extracted in the Excel sheet, in
the copy of Expt2\_final.xlsx, under the Workload, SA, FullData, and
FullData2 sheets.

\subsection{Data preparation and analysis
plan}\label{data-preparation-and-analysis-plan}

In R, data were read and cleaned using FullData2 had many columns that
were imported as character, and they were coerced to numeric (and NA
coping where coercion was needed) De Jonge and Van Der Loo, 2013). In
the case of factors, the participant IDs were transformed to them.
Analyses involved descriptive statistics, visualisations, one-way ANOVAs
(condition as factor), Tukey post-hoc, Levene and ShapiroWilk test of
assumptions, Bayesian ANOVA (BayesFactor package) of workload,
individual ANOVA of NASA-TLX subscales, polychoric correlation of SA
items (attempted), analysis of reliability of NASA-TLX subscales
(Cronbachs alpha) and Pearson correlation of key continuous variables
(score, NASA-TLX, SART, choice accuracy, response time
(Juarros-Basterretxea et al., 2024). The tests of statistics were all
performed with α =.05. In the case of Bayesian Interpretation, standard
interpretation (e.g., BF10, which is greater than 3, is moderate
evidence in favor of H1) was made (Wulff and Taylor, 2024).

\section{Results}\label{results}

All descriptive statistics and inferential findings below are directly
taken off the R analysis results as outlined above. \#\# Descriptive
statistics Descriptive statistics of total workload (as the NASA-TLX
total score) showed that there was a variation between conditions. The
highest workload was recorded in the Extreme\_Hare condition ( E; M =
46.2, SD = 15.1, n = 10), then the Extreme\_Stag condition ( R; M =
34.2, SD = 11.0, n = 10) and the Extreme condition ( O; M = 38.1, SD =
14.8, n = 10). Analysis of NASA-TLX subscale indicated that mental
demand was greatest in the E condition (M= 57.0), a bit lesser in O (M=
54.0), and the least in R (M= 48.0). A similar pattern was applied to
effort, where E = 59, O = 52 and R =34. Conditional selective summary of
these subscale means is as shown below.

The situation awareness, as assessed by way of SART total score,
revealed the opposite situation, with the lowest scores in E (M = 10.2,
SD = 4.24, n = 10) and with the highest scores in O (M = 15.1, SD =
2.88, n = 10) and R (M = 14.6, SD = 3.92, n = 10). According to
FullData, the Scores in the areas of performance were lowest in
Extreme\_Hare (E: M = 691.2, SD = 88.8), higher in the category of
Ambiguous (O: M = 752.0, SD = 95.7), and largest in the category of
Extreme Stag (R: M = 795.6, SD = 36.0). The means of FullData2 were
found to support these individual-participant means (E = 691.0; O =
752.0; R = 796.0).

\begin{Shaded}
\begin{Highlighting}[]
\CommentTok{\# ==============================}
\CommentTok{\# 1. Load \& Install Packages}
\CommentTok{\# ==============================}
\FunctionTok{library}\NormalTok{(tidyverse)    }\CommentTok{\# Data manipulation \& ggplot2}
\end{Highlighting}
\end{Shaded}

\begin{verbatim}
## Warning: package 'tidyverse' was built under R version 4.3.3
\end{verbatim}

\begin{verbatim}
## Warning: package 'tidyr' was built under R version 4.3.3
\end{verbatim}

\begin{verbatim}
## Warning: package 'readr' was built under R version 4.3.3
\end{verbatim}

\begin{verbatim}
## Warning: package 'dplyr' was built under R version 4.3.3
\end{verbatim}

\begin{verbatim}
## Warning: package 'stringr' was built under R version 4.3.3
\end{verbatim}

\begin{verbatim}
## -- Attaching core tidyverse packages ------------------------ tidyverse 2.0.0 --
## v dplyr     1.1.4     v readr     2.1.5
## v forcats   1.0.0     v stringr   1.5.1
## v ggplot2   4.0.0     v tibble    3.2.1
## v lubridate 1.9.3     v tidyr     1.3.1
## v purrr     1.0.2     
## -- Conflicts ------------------------------------------ tidyverse_conflicts() --
## x dplyr::filter() masks stats::filter()
## x dplyr::lag()    masks stats::lag()
## i Use the conflicted package (<http://conflicted.r-lib.org/>) to force all conflicts to become errors
\end{verbatim}

\begin{Shaded}
\begin{Highlighting}[]
\FunctionTok{library}\NormalTok{(janitor)      }\CommentTok{\# Cleaning column names}
\end{Highlighting}
\end{Shaded}

\begin{verbatim}
## Warning: package 'janitor' was built under R version 4.3.3
\end{verbatim}

\begin{verbatim}
## 
## Attaching package: 'janitor'
## 
## The following objects are masked from 'package:stats':
## 
##     chisq.test, fisher.test
\end{verbatim}

\begin{Shaded}
\begin{Highlighting}[]
\FunctionTok{library}\NormalTok{(skimr)        }\CommentTok{\# Overview of data}
\FunctionTok{library}\NormalTok{(lme4)         }\CommentTok{\# Linear mixed{-}effects models}
\end{Highlighting}
\end{Shaded}

\begin{verbatim}
## Warning: package 'lme4' was built under R version 4.3.3
\end{verbatim}

\begin{verbatim}
## Loading required package: Matrix
## 
## Attaching package: 'Matrix'
## 
## The following objects are masked from 'package:tidyr':
## 
##     expand, pack, unpack
\end{verbatim}

\begin{verbatim}
## Warning in check_dep_version(): ABI version mismatch: 
## lme4 was built with Matrix ABI version 1
## Current Matrix ABI version is 0
## Please re-install lme4 from source or restore original 'Matrix' package
\end{verbatim}

\begin{Shaded}
\begin{Highlighting}[]
\FunctionTok{library}\NormalTok{(lmerTest)     }\CommentTok{\# p{-}values for lme4 models}
\end{Highlighting}
\end{Shaded}

\begin{verbatim}
## Warning: package 'lmerTest' was built under R version 4.3.3
\end{verbatim}

\begin{verbatim}
## 
## Attaching package: 'lmerTest'
## 
## The following object is masked from 'package:lme4':
## 
##     lmer
## 
## The following object is masked from 'package:stats':
## 
##     step
\end{verbatim}

\begin{Shaded}
\begin{Highlighting}[]
\FunctionTok{library}\NormalTok{(BayesFactor)  }\CommentTok{\# Bayesian analysis}
\end{Highlighting}
\end{Shaded}

\begin{verbatim}
## Warning: package 'BayesFactor' was built under R version 4.3.3
\end{verbatim}

\begin{verbatim}
## Loading required package: coda
\end{verbatim}

\begin{verbatim}
## Warning: package 'coda' was built under R version 4.3.3
\end{verbatim}

\begin{verbatim}
## Warning in .recacheSubclasses(def@className, def, env): undefined subclass
## "ndiMatrix" of class "replValueSp"; definition not updated
\end{verbatim}

\begin{verbatim}
## ************
## Welcome to BayesFactor 0.9.12-4.7. If you have questions, please contact Richard Morey (richarddmorey@gmail.com).
## 
## Type BFManual() to open the manual.
## ************
\end{verbatim}

\begin{Shaded}
\begin{Highlighting}[]
\FunctionTok{library}\NormalTok{(psych)        }\CommentTok{\# Tetrachoric correlation}
\end{Highlighting}
\end{Shaded}

\begin{verbatim}
## Warning: package 'psych' was built under R version 4.3.3
\end{verbatim}

\begin{verbatim}
## 
## Attaching package: 'psych'
## 
## The following objects are masked from 'package:ggplot2':
## 
##     %+%, alpha
\end{verbatim}

\begin{Shaded}
\begin{Highlighting}[]
\FunctionTok{library}\NormalTok{(polycor)      }\CommentTok{\# Polychoric correlation}
\end{Highlighting}
\end{Shaded}

\begin{verbatim}
## 
## Attaching package: 'polycor'
## 
## The following object is masked from 'package:psych':
## 
##     polyserial
\end{verbatim}

\begin{Shaded}
\begin{Highlighting}[]
\FunctionTok{library}\NormalTok{(reshape2)     }\CommentTok{\# Data reshaping for plotting}
\end{Highlighting}
\end{Shaded}

\begin{verbatim}
## Warning: package 'reshape2' was built under R version 4.3.3
\end{verbatim}

\begin{verbatim}
## 
## Attaching package: 'reshape2'
## 
## The following object is masked from 'package:tidyr':
## 
##     smiths
\end{verbatim}

\begin{Shaded}
\begin{Highlighting}[]
\FunctionTok{library}\NormalTok{(readxl)}
\end{Highlighting}
\end{Shaded}

\begin{verbatim}
## Warning: package 'readxl' was built under R version 4.3.3
\end{verbatim}

\begin{Shaded}
\begin{Highlighting}[]
\FunctionTok{library}\NormalTok{(kableExtra)}
\end{Highlighting}
\end{Shaded}

\begin{verbatim}
## Warning: package 'kableExtra' was built under R version 4.3.3
\end{verbatim}

\begin{verbatim}
## 
## Attaching package: 'kableExtra'
## 
## The following object is masked from 'package:dplyr':
## 
##     group_rows
\end{verbatim}

\begin{Shaded}
\begin{Highlighting}[]
\CommentTok{\# ==============================}
\CommentTok{\# 2. Load Excel Data}
\CommentTok{\# ==============================}
\NormalTok{file\_path }\OtherTok{\textless{}{-}} \StringTok{"C:/Users/brian/Downloads/Copy of Expt2\_final.xlsx"}

\NormalTok{workload  }\OtherTok{\textless{}{-}} \FunctionTok{read\_excel}\NormalTok{(file\_path, }\AttributeTok{sheet =} \StringTok{"Workload"}\NormalTok{)  }\SpecialCharTok{\%\textgreater{}\%} \FunctionTok{clean\_names}\NormalTok{()}
\NormalTok{sa        }\OtherTok{\textless{}{-}} \FunctionTok{read\_excel}\NormalTok{(file\_path, }\AttributeTok{sheet =} \StringTok{"SA"}\NormalTok{)        }\SpecialCharTok{\%\textgreater{}\%} \FunctionTok{clean\_names}\NormalTok{()}
\NormalTok{fulldata  }\OtherTok{\textless{}{-}} \FunctionTok{read\_excel}\NormalTok{(file\_path, }\AttributeTok{sheet =} \StringTok{"FullData"}\NormalTok{)  }\SpecialCharTok{\%\textgreater{}\%} \FunctionTok{clean\_names}\NormalTok{()}
\NormalTok{fulldata2 }\OtherTok{\textless{}{-}} \FunctionTok{read\_excel}\NormalTok{(file\_path, }\AttributeTok{sheet =} \StringTok{"FullData2"}\NormalTok{) }\SpecialCharTok{\%\textgreater{}\%} \FunctionTok{clean\_names}\NormalTok{()}

\NormalTok{fulldata2 }\OtherTok{\textless{}{-}}\NormalTok{ fulldata2 }\SpecialCharTok{\%\textgreater{}\%}
  \FunctionTok{mutate}\NormalTok{(}\FunctionTok{across}\NormalTok{(}\FunctionTok{c}\NormalTok{(score, team\_selection\_clicks, choice\_accuracy, time\_respond\_env,}
\NormalTok{                  proportion\_stag, proportion\_hare, proportion\_dis\_adv,}
\NormalTok{                  nasa\_tlx, sart), as.numeric))}
\end{Highlighting}
\end{Shaded}

\begin{verbatim}
## Warning: There were 9 warnings in `mutate()`.
## The first warning was:
## i In argument: `across(...)`.
## Caused by warning:
## ! NAs introduced by coercion
## i Run `dplyr::last_dplyr_warnings()` to see the 8 remaining warnings.
\end{verbatim}

\begin{Shaded}
\begin{Highlighting}[]
\FunctionTok{cat}\NormalTok{(}\StringTok{"Data loaded successfully!}\SpecialCharTok{\textbackslash{}n}\StringTok{"}\NormalTok{)}
\end{Highlighting}
\end{Shaded}

Data loaded successfully!

\begin{Shaded}
\begin{Highlighting}[]
\CommentTok{\# ==============================}
\CommentTok{\# 3. NASA{-}TLX Table by Condition}
\CommentTok{\# ==============================}
\NormalTok{tlx\_table }\OtherTok{\textless{}{-}}\NormalTok{ workload }\SpecialCharTok{\%\textgreater{}\%}
  \FunctionTok{select}\NormalTok{(condition, mental\_demand, effort, physical\_demand) }\SpecialCharTok{\%\textgreater{}\%}
  \FunctionTok{group\_by}\NormalTok{(condition) }\SpecialCharTok{\%\textgreater{}\%}
  \FunctionTok{summarise}\NormalTok{(}
    \AttributeTok{Mental\_Demand =} \FunctionTok{round}\NormalTok{(}\FunctionTok{mean}\NormalTok{(mental\_demand, }\AttributeTok{na.rm =} \ConstantTok{TRUE}\NormalTok{), }\DecValTok{1}\NormalTok{),}
    \AttributeTok{Effort        =} \FunctionTok{round}\NormalTok{(}\FunctionTok{mean}\NormalTok{(effort, }\AttributeTok{na.rm =} \ConstantTok{TRUE}\NormalTok{), }\DecValTok{1}\NormalTok{),}
    \AttributeTok{Physical\_Demand =} \FunctionTok{round}\NormalTok{(}\FunctionTok{mean}\NormalTok{(physical\_demand, }\AttributeTok{na.rm =} \ConstantTok{TRUE}\NormalTok{), }\DecValTok{1}\NormalTok{),}
    \AttributeTok{.groups =} \StringTok{\textquotesingle{}drop\textquotesingle{}}
\NormalTok{  ) }\SpecialCharTok{\%\textgreater{}\%}
  \FunctionTok{mutate}\NormalTok{(}\AttributeTok{condition =} \FunctionTok{case\_when}\NormalTok{(}
\NormalTok{    condition }\SpecialCharTok{==} \StringTok{"Extreme\_Hare"}  \SpecialCharTok{\textasciitilde{}} \StringTok{"E"}\NormalTok{,}
\NormalTok{    condition }\SpecialCharTok{==} \StringTok{"Ambiguous"}     \SpecialCharTok{\textasciitilde{}} \StringTok{"O"}\NormalTok{,}
\NormalTok{    condition }\SpecialCharTok{==} \StringTok{"Extreme\_Stag"}  \SpecialCharTok{\textasciitilde{}} \StringTok{"R"}\NormalTok{,}
    \ConstantTok{TRUE} \SpecialCharTok{\textasciitilde{}}\NormalTok{ condition}
\NormalTok{  )) }\SpecialCharTok{\%\textgreater{}\%}
  \FunctionTok{pivot\_longer}\NormalTok{(}\AttributeTok{cols =} \SpecialCharTok{{-}}\NormalTok{condition, }\AttributeTok{names\_to =} \StringTok{"Subscale"}\NormalTok{, }\AttributeTok{values\_to =} \StringTok{"Mean"}\NormalTok{) }\SpecialCharTok{\%\textgreater{}\%}
  \FunctionTok{pivot\_wider}\NormalTok{(}\AttributeTok{names\_from =}\NormalTok{ condition, }\AttributeTok{values\_from =}\NormalTok{ Mean)}

\NormalTok{tlx\_table }\SpecialCharTok{\%\textgreater{}\%}
  \FunctionTok{kable}\NormalTok{(}\AttributeTok{caption =} \StringTok{"Table 1: NASA{-}TLX Subscale Means by Condition"}\NormalTok{, }\AttributeTok{booktabs =} \ConstantTok{TRUE}\NormalTok{) }\SpecialCharTok{\%\textgreater{}\%}
  \FunctionTok{kable\_styling}\NormalTok{(}\AttributeTok{full\_width =} \ConstantTok{FALSE}\NormalTok{, }\AttributeTok{position =} \StringTok{"center"}\NormalTok{)}
\end{Highlighting}
\end{Shaded}

\begin{longtable}[t]{lrrr}
\caption{\label{tab:full-analysis}Table 1: NASA-TLX Subscale Means by Condition}\\
\toprule
Subscale & E & O & R\\
\midrule
Mental\_Demand & 57.0 & 54.0 & 48\\
Effort & 59.0 & 52.0 & 34\\
Physical\_Demand & 28.5 & 16.5 & 13\\
\bottomrule
\end{longtable}

\begin{Shaded}
\begin{Highlighting}[]
\CommentTok{\# ==============================}
\CommentTok{\# 4. Descriptive Statistics}
\CommentTok{\# ==============================}
\NormalTok{workload\_summary }\OtherTok{\textless{}{-}}\NormalTok{ workload }\SpecialCharTok{\%\textgreater{}\%}
  \FunctionTok{group\_by}\NormalTok{(condition) }\SpecialCharTok{\%\textgreater{}\%}
  \FunctionTok{summarise}\NormalTok{(}
    \AttributeTok{n =} \FunctionTok{n}\NormalTok{(),}
    \AttributeTok{mean\_total\_score =} \FunctionTok{mean}\NormalTok{(total\_score\_sum\_6, }\AttributeTok{na.rm =} \ConstantTok{TRUE}\NormalTok{),}
    \AttributeTok{sd\_total\_score =} \FunctionTok{sd}\NormalTok{(total\_score\_sum\_6, }\AttributeTok{na.rm =} \ConstantTok{TRUE}\NormalTok{),}
    \AttributeTok{se\_total\_score =}\NormalTok{ sd\_total\_score }\SpecialCharTok{/} \FunctionTok{sqrt}\NormalTok{(n),}
    \AttributeTok{.groups =} \StringTok{\textquotesingle{}drop\textquotesingle{}}
\NormalTok{  )}

\NormalTok{fulldata\_summary }\OtherTok{\textless{}{-}}\NormalTok{ fulldata }\SpecialCharTok{\%\textgreater{}\%}
  \FunctionTok{group\_by}\NormalTok{(condition) }\SpecialCharTok{\%\textgreater{}\%}
  \FunctionTok{summarise}\NormalTok{(}
    \AttributeTok{n =} \FunctionTok{n}\NormalTok{(),}
    \AttributeTok{mean\_score =} \FunctionTok{mean}\NormalTok{(score, }\AttributeTok{na.rm =} \ConstantTok{TRUE}\NormalTok{),}
    \AttributeTok{sd\_score =} \FunctionTok{sd}\NormalTok{(score, }\AttributeTok{na.rm =} \ConstantTok{TRUE}\NormalTok{),}
    \AttributeTok{se\_score =}\NormalTok{ sd\_score }\SpecialCharTok{/} \FunctionTok{sqrt}\NormalTok{(n),}
    \AttributeTok{.groups =} \StringTok{\textquotesingle{}drop\textquotesingle{}}
\NormalTok{  )}

\NormalTok{sa\_summary }\OtherTok{\textless{}{-}}\NormalTok{ sa }\SpecialCharTok{\%\textgreater{}\%}
  \FunctionTok{group\_by}\NormalTok{(demand\_condition) }\SpecialCharTok{\%\textgreater{}\%}
  \FunctionTok{summarise}\NormalTok{(}
    \AttributeTok{n =} \FunctionTok{n}\NormalTok{(),}
    \AttributeTok{mean\_sa\_total =} \FunctionTok{mean}\NormalTok{(sa\_total, }\AttributeTok{na.rm =} \ConstantTok{TRUE}\NormalTok{),}
    \AttributeTok{sd\_sa\_total =} \FunctionTok{sd}\NormalTok{(sa\_total, }\AttributeTok{na.rm =} \ConstantTok{TRUE}\NormalTok{),}
    \AttributeTok{se\_sa\_total =}\NormalTok{ sd\_sa\_total }\SpecialCharTok{/} \FunctionTok{sqrt}\NormalTok{(n),}
    \AttributeTok{.groups =} \StringTok{\textquotesingle{}drop\textquotesingle{}}
\NormalTok{  )}

\FunctionTok{cat}\NormalTok{(}\StringTok{"Descriptive statistics computed.}\SpecialCharTok{\textbackslash{}n}\StringTok{"}\NormalTok{)}
\end{Highlighting}
\end{Shaded}

Descriptive statistics computed.

\begin{Shaded}
\begin{Highlighting}[]
\CommentTok{\# ==============================}
\CommentTok{\# 5. Plots}
\CommentTok{\# ==============================}
\CommentTok{\# Plot 1: Workload Total Score}
\NormalTok{p1 }\OtherTok{\textless{}{-}} \FunctionTok{ggplot}\NormalTok{(workload, }\FunctionTok{aes}\NormalTok{(}\AttributeTok{x =}\NormalTok{ condition, }\AttributeTok{y =}\NormalTok{ total\_score\_sum\_6, }\AttributeTok{fill =}\NormalTok{ condition)) }\SpecialCharTok{+}
  \FunctionTok{geom\_boxplot}\NormalTok{(}\AttributeTok{alpha =} \FloatTok{0.7}\NormalTok{) }\SpecialCharTok{+}
  \FunctionTok{geom\_jitter}\NormalTok{(}\AttributeTok{width =} \FloatTok{0.2}\NormalTok{, }\AttributeTok{alpha =} \FloatTok{0.3}\NormalTok{) }\SpecialCharTok{+}
  \FunctionTok{theme\_minimal}\NormalTok{(}\AttributeTok{base\_size =} \DecValTok{12}\NormalTok{) }\SpecialCharTok{+}
  \FunctionTok{labs}\NormalTok{(}\AttributeTok{title =} \StringTok{"Workload Total Score by Condition"}\NormalTok{, }\AttributeTok{x =} \StringTok{"Condition"}\NormalTok{, }\AttributeTok{y =} \StringTok{"Total Score (sum/6)"}\NormalTok{) }\SpecialCharTok{+}
  \FunctionTok{theme}\NormalTok{(}\AttributeTok{legend.position =} \StringTok{"none"}\NormalTok{)}
\FunctionTok{print}\NormalTok{(p1)}
\end{Highlighting}
\end{Shaded}

\includegraphics{YY56_files/figure-latex/full-analysis-1.pdf}

\begin{Shaded}
\begin{Highlighting}[]
\CommentTok{\# Plot 2: NASA{-}TLX Subscales}
\NormalTok{workload\_melt }\OtherTok{\textless{}{-}}\NormalTok{ workload }\SpecialCharTok{\%\textgreater{}\%}
  \FunctionTok{select}\NormalTok{(participant, condition, mental\_demand, physical\_demand, temporal\_demand,}
\NormalTok{         performance, effort, frustration) }\SpecialCharTok{\%\textgreater{}\%}
  \FunctionTok{pivot\_longer}\NormalTok{(}\AttributeTok{cols =}\NormalTok{ mental\_demand}\SpecialCharTok{:}\NormalTok{frustration, }\AttributeTok{names\_to =} \StringTok{"subscale"}\NormalTok{, }\AttributeTok{values\_to =} \StringTok{"value"}\NormalTok{)}

\NormalTok{p2 }\OtherTok{\textless{}{-}} \FunctionTok{ggplot}\NormalTok{(workload\_melt, }\FunctionTok{aes}\NormalTok{(}\AttributeTok{x =}\NormalTok{ subscale, }\AttributeTok{y =}\NormalTok{ value, }\AttributeTok{fill =}\NormalTok{ condition)) }\SpecialCharTok{+}
  \FunctionTok{stat\_summary}\NormalTok{(}\AttributeTok{fun =}\NormalTok{ mean, }\AttributeTok{geom =} \StringTok{"bar"}\NormalTok{, }\AttributeTok{position =} \FunctionTok{position\_dodge}\NormalTok{(}\FloatTok{0.8}\NormalTok{), }\AttributeTok{alpha =} \FloatTok{0.8}\NormalTok{) }\SpecialCharTok{+}
  \FunctionTok{stat\_summary}\NormalTok{(}\AttributeTok{fun.data =}\NormalTok{ mean\_se, }\AttributeTok{geom =} \StringTok{"errorbar"}\NormalTok{, }\AttributeTok{width =} \FloatTok{0.2}\NormalTok{, }\AttributeTok{position =} \FunctionTok{position\_dodge}\NormalTok{(}\FloatTok{0.8}\NormalTok{)) }\SpecialCharTok{+}
  \FunctionTok{theme\_minimal}\NormalTok{(}\AttributeTok{base\_size =} \DecValTok{12}\NormalTok{) }\SpecialCharTok{+}
  \FunctionTok{labs}\NormalTok{(}\AttributeTok{title =} \StringTok{"Mean NASA{-}TLX Subscale Scores by Condition"}\NormalTok{, }\AttributeTok{x =} \StringTok{"Subscale"}\NormalTok{, }\AttributeTok{y =} \StringTok{"Mean Score"}\NormalTok{) }\SpecialCharTok{+}
  \FunctionTok{theme}\NormalTok{(}\AttributeTok{axis.text.x =} \FunctionTok{element\_text}\NormalTok{(}\AttributeTok{angle =} \DecValTok{45}\NormalTok{, }\AttributeTok{hjust =} \DecValTok{1}\NormalTok{))}
\FunctionTok{print}\NormalTok{(p2)}
\end{Highlighting}
\end{Shaded}

\includegraphics{YY56_files/figure-latex/full-analysis-2.pdf}

\begin{Shaded}
\begin{Highlighting}[]
\CommentTok{\# Plot 3: Performance Scores}
\NormalTok{p3 }\OtherTok{\textless{}{-}} \FunctionTok{ggplot}\NormalTok{(fulldata, }\FunctionTok{aes}\NormalTok{(}\AttributeTok{x =}\NormalTok{ condition, }\AttributeTok{y =}\NormalTok{ score, }\AttributeTok{fill =}\NormalTok{ condition)) }\SpecialCharTok{+}
  \FunctionTok{geom\_violin}\NormalTok{(}\AttributeTok{alpha =} \FloatTok{0.6}\NormalTok{) }\SpecialCharTok{+}
  \FunctionTok{geom\_boxplot}\NormalTok{(}\AttributeTok{width =} \FloatTok{0.2}\NormalTok{, }\AttributeTok{alpha =} \FloatTok{0.8}\NormalTok{) }\SpecialCharTok{+}
  \FunctionTok{theme\_minimal}\NormalTok{(}\AttributeTok{base\_size =} \DecValTok{12}\NormalTok{) }\SpecialCharTok{+}
  \FunctionTok{labs}\NormalTok{(}\AttributeTok{title =} \StringTok{"Performance Scores by Condition"}\NormalTok{, }\AttributeTok{x =} \StringTok{"Condition"}\NormalTok{, }\AttributeTok{y =} \StringTok{"Score"}\NormalTok{) }\SpecialCharTok{+}
  \FunctionTok{theme}\NormalTok{(}\AttributeTok{legend.position =} \StringTok{"none"}\NormalTok{)}
\FunctionTok{print}\NormalTok{(p3)}
\end{Highlighting}
\end{Shaded}

\includegraphics{YY56_files/figure-latex/full-analysis-3.pdf}

\begin{Shaded}
\begin{Highlighting}[]
\CommentTok{\# Plot 4: Correlation Matrix}
\NormalTok{cor\_data }\OtherTok{\textless{}{-}}\NormalTok{ fulldata }\SpecialCharTok{\%\textgreater{}\%} \FunctionTok{select}\NormalTok{(score, nasa\_tlx, sart, choice\_accuracy, time\_respond\_env)}
\NormalTok{cor\_matrix }\OtherTok{\textless{}{-}} \FunctionTok{cor}\NormalTok{(cor\_data, }\AttributeTok{use =} \StringTok{"pairwise.complete.obs"}\NormalTok{)}
\NormalTok{cor\_melt }\OtherTok{\textless{}{-}} \FunctionTok{melt}\NormalTok{(cor\_matrix)}

\NormalTok{p4 }\OtherTok{\textless{}{-}} \FunctionTok{ggplot}\NormalTok{(cor\_melt, }\FunctionTok{aes}\NormalTok{(}\AttributeTok{x =}\NormalTok{ Var1, }\AttributeTok{y =}\NormalTok{ Var2, }\AttributeTok{fill =}\NormalTok{ value)) }\SpecialCharTok{+}
  \FunctionTok{geom\_tile}\NormalTok{() }\SpecialCharTok{+}
  \FunctionTok{geom\_text}\NormalTok{(}\FunctionTok{aes}\NormalTok{(}\AttributeTok{label =} \FunctionTok{round}\NormalTok{(value, }\DecValTok{2}\NormalTok{)), }\AttributeTok{color =} \StringTok{"black"}\NormalTok{, }\AttributeTok{size =} \DecValTok{3}\NormalTok{) }\SpecialCharTok{+}
  \FunctionTok{scale\_fill\_gradient2}\NormalTok{(}\AttributeTok{low =} \StringTok{"blue"}\NormalTok{, }\AttributeTok{mid =} \StringTok{"white"}\NormalTok{, }\AttributeTok{high =} \StringTok{"red"}\NormalTok{, }\AttributeTok{midpoint =} \DecValTok{0}\NormalTok{, }\AttributeTok{limit =} \FunctionTok{c}\NormalTok{(}\SpecialCharTok{{-}}\DecValTok{1}\NormalTok{,}\DecValTok{1}\NormalTok{)) }\SpecialCharTok{+}
  \FunctionTok{theme\_minimal}\NormalTok{() }\SpecialCharTok{+}
  \FunctionTok{labs}\NormalTok{(}\AttributeTok{title =} \StringTok{"Correlation Matrix: Performance Measures"}\NormalTok{) }\SpecialCharTok{+}
  \FunctionTok{theme}\NormalTok{(}\AttributeTok{axis.text.x =} \FunctionTok{element\_text}\NormalTok{(}\AttributeTok{angle =} \DecValTok{45}\NormalTok{, }\AttributeTok{hjust =} \DecValTok{1}\NormalTok{))}
\FunctionTok{print}\NormalTok{(p4)}
\end{Highlighting}
\end{Shaded}

\includegraphics{YY56_files/figure-latex/full-analysis-4.pdf}

\begin{Shaded}
\begin{Highlighting}[]
\CommentTok{\# ==============================}
\CommentTok{\# 6. Statistical Analyses}
\CommentTok{\# ==============================}
\CommentTok{\# Convert factors}
\NormalTok{workload}\SpecialCharTok{$}\NormalTok{condition }\OtherTok{\textless{}{-}} \FunctionTok{as.factor}\NormalTok{(workload}\SpecialCharTok{$}\NormalTok{condition)}
\NormalTok{sa}\SpecialCharTok{$}\NormalTok{demand\_condition }\OtherTok{\textless{}{-}} \FunctionTok{as.factor}\NormalTok{(sa}\SpecialCharTok{$}\NormalTok{demand\_condition)}
\NormalTok{fulldata}\SpecialCharTok{$}\NormalTok{condition }\OtherTok{\textless{}{-}} \FunctionTok{as.factor}\NormalTok{(fulldata}\SpecialCharTok{$}\NormalTok{condition)}

\CommentTok{\# ANOVA: Workload Total Score}
\FunctionTok{cat}\NormalTok{(}\StringTok{"}\SpecialCharTok{\textbackslash{}n}\StringTok{{-}{-}{-} ANOVA: Workload Total Score {-}{-}{-}}\SpecialCharTok{\textbackslash{}n}\StringTok{"}\NormalTok{)}
\end{Highlighting}
\end{Shaded}

--- ANOVA: Workload Total Score ---

\begin{Shaded}
\begin{Highlighting}[]
\NormalTok{aov\_workload }\OtherTok{\textless{}{-}} \FunctionTok{aov}\NormalTok{(total\_score\_sum\_6 }\SpecialCharTok{\textasciitilde{}}\NormalTok{ condition, }\AttributeTok{data =}\NormalTok{ workload)}
\FunctionTok{print}\NormalTok{(}\FunctionTok{summary}\NormalTok{(aov\_workload))}
\end{Highlighting}
\end{Shaded}

\begin{verbatim}
        Df Sum Sq Mean Sq F value Pr(>F)
\end{verbatim}

condition 2 747 373.6 1.976 0.158 Residuals 27 5106 189.1

\begin{Shaded}
\begin{Highlighting}[]
\CommentTok{\# Post{-}hoc Tukey}
\NormalTok{tukey\_workload }\OtherTok{\textless{}{-}} \FunctionTok{TukeyHSD}\NormalTok{(aov\_workload)}
\FunctionTok{print}\NormalTok{(tukey\_workload)}
\end{Highlighting}
\end{Shaded}

Tukey multiple comparisons of means 95\% family-wise confidence level

Fit: aov(formula = total\_score\_sum\_6 \textasciitilde{} condition,
data = workload)

\$condition diff lwr upr p adj O-E -8.083333 -23.33159 7.164922
0.3996422 R-E -11.983333 -27.23159 3.264922 0.1446954 R-O -3.900000
-19.14825 11.348255 0.8027425

\begin{Shaded}
\begin{Highlighting}[]
\CommentTok{\# Effect size}
\ControlFlowTok{if}\NormalTok{(}\FunctionTok{require}\NormalTok{(effectsize, }\AttributeTok{quietly =} \ConstantTok{TRUE}\NormalTok{)) \{}
\NormalTok{  eta\_sq }\OtherTok{\textless{}{-}} \FunctionTok{eta\_squared}\NormalTok{(aov\_workload)}
  \FunctionTok{print}\NormalTok{(eta\_sq)}
\NormalTok{\}}
\end{Highlighting}
\end{Shaded}

\begin{verbatim}
## 
## Attaching package: 'effectsize'
## 
## The following object is masked from 'package:psych':
## 
##     phi
## 
## For one-way between subjects designs, partial eta squared is equivalent
##   to eta squared. Returning eta squared.
\end{verbatim}

\section{Effect Size for ANOVA}\label{effect-size-for-anova}

\subsection{Parameter \textbar{} Eta2 \textbar{} 95\%
CI}\label{parameter-eta2-95-ci}

condition \textbar{} 0.13 \textbar{} {[}0.00, 1.00{]}

\begin{itemize}
\tightlist
\item
  One-sided CIs: upper bound fixed at {[}1.00{]}.
\end{itemize}

\begin{Shaded}
\begin{Highlighting}[]
\CommentTok{\# ANOVA: Situation Awareness}
\FunctionTok{cat}\NormalTok{(}\StringTok{"}\SpecialCharTok{\textbackslash{}n}\StringTok{{-}{-}{-} ANOVA: Situation Awareness {-}{-}{-}}\SpecialCharTok{\textbackslash{}n}\StringTok{"}\NormalTok{)}
\end{Highlighting}
\end{Shaded}

--- ANOVA: Situation Awareness ---

\begin{Shaded}
\begin{Highlighting}[]
\NormalTok{aov\_sa }\OtherTok{\textless{}{-}} \FunctionTok{aov}\NormalTok{(sa\_total }\SpecialCharTok{\textasciitilde{}}\NormalTok{ demand\_condition, }\AttributeTok{data =}\NormalTok{ sa)}
\FunctionTok{print}\NormalTok{(}\FunctionTok{summary}\NormalTok{(aov\_sa))}
\end{Highlighting}
\end{Shaded}

\begin{verbatim}
             Df Sum Sq Mean Sq F value Pr(>F)  
\end{verbatim}

demand\_condition 2 145.4 72.70 5.236 0.012 * Residuals 27 374.9 13.89\\
--- Signif. codes: 0 `\emph{\textbf{' 0.001 '}' 0.01 '}' 0.05 `.' 0.1 '
' 1

\begin{Shaded}
\begin{Highlighting}[]
\NormalTok{tukey\_sa }\OtherTok{\textless{}{-}} \FunctionTok{TukeyHSD}\NormalTok{(aov\_sa)}
\FunctionTok{print}\NormalTok{(tukey\_sa)}
\end{Highlighting}
\end{Shaded}

Tukey multiple comparisons of means 95\% family-wise confidence level

Fit: aov(formula = sa\_total \textasciitilde{} demand\_condition, data =
sa)

\$demand\_condition diff lwr upr p adj O-E 4.9 0.7681882 9.031812
0.0176728 R-E 4.4 0.2681882 8.531812 0.0351504 R-O -0.5 -4.6318118
3.631812 0.9516749

\begin{Shaded}
\begin{Highlighting}[]
\CommentTok{\# ANOVA: Performance}
\FunctionTok{cat}\NormalTok{(}\StringTok{"}\SpecialCharTok{\textbackslash{}n}\StringTok{{-}{-}{-} ANOVA: Performance Scores {-}{-}{-}}\SpecialCharTok{\textbackslash{}n}\StringTok{"}\NormalTok{)}
\end{Highlighting}
\end{Shaded}

--- ANOVA: Performance Scores ---

\begin{Shaded}
\begin{Highlighting}[]
\NormalTok{aov\_perf }\OtherTok{\textless{}{-}} \FunctionTok{aov}\NormalTok{(score }\SpecialCharTok{\textasciitilde{}}\NormalTok{ condition, }\AttributeTok{data =}\NormalTok{ fulldata)}
\FunctionTok{print}\NormalTok{(}\FunctionTok{summary}\NormalTok{(aov\_perf))}
\end{Highlighting}
\end{Shaded}

\begin{verbatim}
        Df Sum Sq Mean Sq F value Pr(>F)  
\end{verbatim}

condition 2 54934 27467 4.492 0.0207 * Residuals 27 165098 6115\\
--- Signif. codes: 0 `\emph{\textbf{' 0.001 '}' 0.01 '}' 0.05 `.' 0.1 '
' 1

\begin{Shaded}
\begin{Highlighting}[]
\NormalTok{tukey\_perf }\OtherTok{\textless{}{-}} \FunctionTok{TukeyHSD}\NormalTok{(aov\_perf)}
\FunctionTok{print}\NormalTok{(tukey\_perf)}
\end{Highlighting}
\end{Shaded}

Tukey multiple comparisons of means 95\% family-wise confidence level

Fit: aov(formula = score \textasciitilde{} condition, data = fulldata)

\$condition diff lwr upr p adj O-E 60.3 -26.40704 147.007 0.2146782 R-E
104.4 17.69296 191.107 0.0158937 R-O 44.1 -42.60704 130.807 0.4288152

\begin{Shaded}
\begin{Highlighting}[]
\FunctionTok{cat}\NormalTok{(}\StringTok{"}\SpecialCharTok{\textbackslash{}n}\StringTok{=== ANALYSIS COMPLETE ===}\SpecialCharTok{\textbackslash{}n}\StringTok{"}\NormalTok{)}
\end{Highlighting}
\end{Shaded}

=== ANALYSIS COMPLETE === The situation awareness, as assessed by way of
SART total score, revealed the opposite situation, with the lowest
scores in E (M = 10.2, SD = 4.24, n = 10) and with the highest scores in
O (M = 15.1, SD = 2.88, n = 10) and R (M = 14.6, SD = 3.92, n = 10).
According to FullData, the Scores in the areas of performance were
lowest in Extreme\_Hare (E: M = 691.2, SD = 88.8), higher in the
category of Ambiguous (O: M = 752.0, SD = 95.7), and largest in the
category of Extreme Stag (R: M = 795.6, SD = 36.0). The means of
FullData2 were found to support these individual-participant means (E =
691.0; O = 752.0; R = 796.0).

\#\#Inferential statistics \#\#\# Workload (NASA-TLX total) One-way
ANOVA of the differences in NASA-TLX total scores at different
conditions showed no significance, F(2,27) = 1.98, p = 0.158.
Homogeneity of variance (Levene, p = 0.774) and approximate normality of
residuals (Shapiro Wilk, W = 0.964, p = 0.388) was checked by assumption
checks. All post-hoc Tukey comparisons were not significant, such as O
-E (diff = -8.08, p = 0.400) and R -E ( diff = -11.98, p = 0.145). Cohen
effect sizes of the workload comparisons were varied with the biggest
value of the d effect being observed between R vs E (d = -0.908) and the
least value observed between Extreme\_Stag vs Extreme Hare (d = 0.393).
Analysis under Bayesian ANOVA gave a Bayes factor of 0.733 ( 0.01\%)
which had weak evidence to support the null model against the condition
model, which was similar to the non-significant ANOVA results. Figure 1
will give the boxplot of the total workload scores across conditions and
indicate the patterns of distribution. The mean score of the
Extreme\_Hare (E) condition is 691.2 (SD = 88.8), and the mean score of
the Ambiguous (O) condition is 752.0 (SD = 95.7), and the latter is
slightly higher, which may be attributed to an anticipated higher
effort, potentially caused by uncertainty in the dynamics of the task.
And the highest mean score which was 796.0 (SD = 36.0) but with a
significantly smaller standard deviation was found in the Extreme\_Stag
(R) condition which implies that the mean workload rating of the
participants was less variable when the environment was evidently biased
towards the high-reward Stag team (Dung and Duong, 2024).

\begin{Shaded}
\begin{Highlighting}[]
\CommentTok{\# ==============================}
\CommentTok{\# 3.2.1 Workload (NASA{-}TLX total)}
\CommentTok{\# ==============================}

\FunctionTok{library}\NormalTok{(car)        }\CommentTok{\# Levene\textquotesingle{}s test}
\FunctionTok{library}\NormalTok{(effectsize)  }\CommentTok{\# Cohen\textquotesingle{}s d}
\FunctionTok{library}\NormalTok{(BayesFactor) }\CommentTok{\# Bayesian ANOVA}
\FunctionTok{library}\NormalTok{(ggplot2)     }\CommentTok{\# Plotting}

\CommentTok{\# Convert condition to factor}
\NormalTok{workload}\SpecialCharTok{$}\NormalTok{condition }\OtherTok{\textless{}{-}} \FunctionTok{as.factor}\NormalTok{(workload}\SpecialCharTok{$}\NormalTok{condition)}

\CommentTok{\# One{-}way ANOVA: Total NASA{-}TLX score}
\NormalTok{aov\_workload }\OtherTok{\textless{}{-}} \FunctionTok{aov}\NormalTok{(total\_score\_sum\_6 }\SpecialCharTok{\textasciitilde{}}\NormalTok{ condition, }\AttributeTok{data =}\NormalTok{ workload)}
\NormalTok{summary\_aov }\OtherTok{\textless{}{-}} \FunctionTok{summary}\NormalTok{(aov\_workload)}
\FunctionTok{print}\NormalTok{(summary\_aov)}
\end{Highlighting}
\end{Shaded}

\begin{verbatim}
##             Df Sum Sq Mean Sq F value Pr(>F)
## condition    2    747   373.6   1.976  0.158
## Residuals   27   5106   189.1
\end{verbatim}

\begin{Shaded}
\begin{Highlighting}[]
\CommentTok{\# Assumption checks}
\CommentTok{\# Levene\textquotesingle{}s test for homogeneity of variance}
\NormalTok{levene\_test }\OtherTok{\textless{}{-}} \FunctionTok{leveneTest}\NormalTok{(total\_score\_sum\_6 }\SpecialCharTok{\textasciitilde{}}\NormalTok{ condition, }\AttributeTok{data =}\NormalTok{ workload)}
\FunctionTok{print}\NormalTok{(levene\_test)}
\end{Highlighting}
\end{Shaded}

\begin{verbatim}
## Levene's Test for Homogeneity of Variance (center = median)
##       Df F value Pr(>F)
## group  2  0.2585 0.7741
##       27
\end{verbatim}

\begin{Shaded}
\begin{Highlighting}[]
\CommentTok{\# Shapiro{-}Wilk test for normality of residuals}
\NormalTok{shapiro\_test }\OtherTok{\textless{}{-}} \FunctionTok{shapiro.test}\NormalTok{(}\FunctionTok{residuals}\NormalTok{(aov\_workload))}
\FunctionTok{cat}\NormalTok{(}\StringTok{"Shapiro{-}Wilk W ="}\NormalTok{, }\FunctionTok{round}\NormalTok{(shapiro\_test}\SpecialCharTok{$}\NormalTok{statistic,}\DecValTok{3}\NormalTok{), }
    \StringTok{", p ="}\NormalTok{, }\FunctionTok{round}\NormalTok{(shapiro\_test}\SpecialCharTok{$}\NormalTok{p.value,}\DecValTok{3}\NormalTok{), }\StringTok{"}\SpecialCharTok{\textbackslash{}n}\StringTok{"}\NormalTok{)}
\end{Highlighting}
\end{Shaded}

\begin{verbatim}
## Shapiro-Wilk W = 0.964 , p = 0.388
\end{verbatim}

\begin{Shaded}
\begin{Highlighting}[]
\CommentTok{\# Post{-}hoc Tukey HSD comparisons}
\NormalTok{tukey\_results }\OtherTok{\textless{}{-}} \FunctionTok{TukeyHSD}\NormalTok{(aov\_workload)}
\FunctionTok{print}\NormalTok{(tukey\_results)}
\end{Highlighting}
\end{Shaded}

\begin{verbatim}
##   Tukey multiple comparisons of means
##     95% family-wise confidence level
## 
## Fit: aov(formula = total_score_sum_6 ~ condition, data = workload)
## 
## $condition
##           diff       lwr       upr     p adj
## O-E  -8.083333 -23.33159  7.164922 0.3996422
## R-E -11.983333 -27.23159  3.264922 0.1446954
## R-O  -3.900000 -19.14825 11.348255 0.8027425
\end{verbatim}

\begin{Shaded}
\begin{Highlighting}[]
\CommentTok{\# Cohen\textquotesingle{}s d effect sizes for pairwise comparisons}
\NormalTok{conditions }\OtherTok{\textless{}{-}} \FunctionTok{levels}\NormalTok{(workload}\SpecialCharTok{$}\NormalTok{condition)}
\FunctionTok{cat}\NormalTok{(}\StringTok{"}\SpecialCharTok{\textbackslash{}n}\StringTok{Cohen\textquotesingle{}s d effect sizes:}\SpecialCharTok{\textbackslash{}n}\StringTok{"}\NormalTok{)}
\end{Highlighting}
\end{Shaded}

\begin{verbatim}
## 
## Cohen's d effect sizes:
\end{verbatim}

\begin{Shaded}
\begin{Highlighting}[]
\ControlFlowTok{for}\NormalTok{ (i }\ControlFlowTok{in} \DecValTok{1}\SpecialCharTok{:}\NormalTok{(}\FunctionTok{length}\NormalTok{(conditions)}\SpecialCharTok{{-}}\DecValTok{1}\NormalTok{)) \{}
  \ControlFlowTok{for}\NormalTok{ (j }\ControlFlowTok{in}\NormalTok{ (i}\SpecialCharTok{+}\DecValTok{1}\NormalTok{)}\SpecialCharTok{:}\FunctionTok{length}\NormalTok{(conditions)) \{}
\NormalTok{    group1 }\OtherTok{\textless{}{-}}\NormalTok{ workload}\SpecialCharTok{$}\NormalTok{total\_score\_sum\_6[workload}\SpecialCharTok{$}\NormalTok{condition }\SpecialCharTok{==}\NormalTok{ conditions[i]]}
\NormalTok{    group2 }\OtherTok{\textless{}{-}}\NormalTok{ workload}\SpecialCharTok{$}\NormalTok{total\_score\_sum\_6[workload}\SpecialCharTok{$}\NormalTok{condition }\SpecialCharTok{==}\NormalTok{ conditions[j]]}
\NormalTok{    d\_val }\OtherTok{\textless{}{-}}\NormalTok{ (}\FunctionTok{mean}\NormalTok{(group2, }\AttributeTok{na.rm=}\ConstantTok{TRUE}\NormalTok{) }\SpecialCharTok{{-}} \FunctionTok{mean}\NormalTok{(group1, }\AttributeTok{na.rm=}\ConstantTok{TRUE}\NormalTok{)) }\SpecialCharTok{/} 
      \FunctionTok{sqrt}\NormalTok{((}\FunctionTok{sd}\NormalTok{(group1, }\AttributeTok{na.rm=}\ConstantTok{TRUE}\NormalTok{)}\SpecialCharTok{\^{}}\DecValTok{2} \SpecialCharTok{+} \FunctionTok{sd}\NormalTok{(group2, }\AttributeTok{na.rm=}\ConstantTok{TRUE}\NormalTok{)}\SpecialCharTok{\^{}}\DecValTok{2}\NormalTok{)}\SpecialCharTok{/}\DecValTok{2}\NormalTok{)}
    \FunctionTok{cat}\NormalTok{(conditions[j], }\StringTok{"vs"}\NormalTok{, conditions[i], }\StringTok{": d ="}\NormalTok{, }\FunctionTok{round}\NormalTok{(d\_val,}\DecValTok{3}\NormalTok{), }\StringTok{"}\SpecialCharTok{\textbackslash{}n}\StringTok{"}\NormalTok{)}
\NormalTok{  \}}
\NormalTok{\}}
\end{Highlighting}
\end{Shaded}

\begin{verbatim}
## O vs E : d = -0.541 
## R vs E : d = -0.908 
## R vs O : d = -0.299
\end{verbatim}

\begin{Shaded}
\begin{Highlighting}[]
\CommentTok{\# Bayesian ANOVA}
\FunctionTok{cat}\NormalTok{(}\StringTok{"}\SpecialCharTok{\textbackslash{}n}\StringTok{Bayesian ANOVA:}\SpecialCharTok{\textbackslash{}n}\StringTok{"}\NormalTok{)}
\end{Highlighting}
\end{Shaded}

\begin{verbatim}
## 
## Bayesian ANOVA:
\end{verbatim}

\begin{Shaded}
\begin{Highlighting}[]
\NormalTok{bf\_workload }\OtherTok{\textless{}{-}} \FunctionTok{tryCatch}\NormalTok{(\{}
  \FunctionTok{anovaBF}\NormalTok{(total\_score\_sum\_6 }\SpecialCharTok{\textasciitilde{}}\NormalTok{ condition, }\AttributeTok{data =}\NormalTok{ workload)}
\NormalTok{\}, }\AttributeTok{error =} \ControlFlowTok{function}\NormalTok{(e) \{}\FunctionTok{cat}\NormalTok{(}\StringTok{"Bayesian ANOVA error:"}\NormalTok{, e}\SpecialCharTok{$}\NormalTok{message, }\StringTok{"}\SpecialCharTok{\textbackslash{}n}\StringTok{"}\NormalTok{)\})}
\FunctionTok{print}\NormalTok{(bf\_workload)}
\end{Highlighting}
\end{Shaded}

\begin{verbatim}
## Bayes factor analysis
## --------------
## [1] condition : 0.732982 ±0.01%
## 
## Against denominator:
##   Intercept only 
## ---
## Bayes factor type: BFlinearModel, JZS
\end{verbatim}

\begin{Shaded}
\begin{Highlighting}[]
\CommentTok{\# ==============================}
\CommentTok{\# Figure 1: Boxplot of NASA{-}TLX Total Scores}
\CommentTok{\# ==============================}
\FunctionTok{ggplot}\NormalTok{(workload, }\FunctionTok{aes}\NormalTok{(}\AttributeTok{x =}\NormalTok{ condition, }\AttributeTok{y =}\NormalTok{ total\_score\_sum\_6, }\AttributeTok{fill =}\NormalTok{ condition)) }\SpecialCharTok{+}
  \FunctionTok{geom\_boxplot}\NormalTok{(}\AttributeTok{alpha =} \FloatTok{0.7}\NormalTok{) }\SpecialCharTok{+}
  \FunctionTok{geom\_jitter}\NormalTok{(}\AttributeTok{width =} \FloatTok{0.2}\NormalTok{, }\AttributeTok{alpha =} \FloatTok{0.3}\NormalTok{) }\SpecialCharTok{+}
  \FunctionTok{theme\_minimal}\NormalTok{(}\AttributeTok{base\_size =} \DecValTok{12}\NormalTok{) }\SpecialCharTok{+}
  \FunctionTok{labs}\NormalTok{(}
    \AttributeTok{title =} \StringTok{"Figure 1: NASA{-}TLX Total Workload Scores by Condition"}\NormalTok{,}
    \AttributeTok{x =} \StringTok{"Condition"}\NormalTok{,}
    \AttributeTok{y =} \StringTok{"Total Workload Score"}
\NormalTok{  ) }\SpecialCharTok{+}
  \FunctionTok{theme}\NormalTok{(}\AttributeTok{legend.position =} \StringTok{"none"}\NormalTok{)}
\end{Highlighting}
\end{Shaded}

\includegraphics{YY56_files/figure-latex/workload-inferential-1.pdf}
\#\#\# Situation Awareness (SART total) The total score of SART was
significantly different in the three environmental conditions, which was
situation awareness. An ANOVA with one independent t-test condition
showed that there was a significant effect of condition, F(2, 27) =
5.236, p=0.012, which means that the awareness of system state of the
participants did depend on whether they managed Extreme\_Hare (E),
Ambiguous (O), or Extreme\_Stag (R) environments. Ambiguous condition
versus Extreme-Hare condition post-hoc Tukey comparisons revealed a
significant mean difference between situation awareness in the situation
of ambiguity (O; M = 15.1, SD = 2.88) and the situation of extreme hare
(E; M = 10.2, SD = 4.24), mean difference = 4.9, 95\% CI {[}0.768,
9.032{]}, p = 0.018. Equally, the SART scores were much higher in the
Extreme\_Stag condition (R; M = 14.6, SD = 3.92) in comparison with E
(mean diff = 4.4, 95\% CI {[}0.268, 8.532{]}, p = 0.035), and the
difference between R and O was not significant (diff = -0.5, p = 0.952).
Homogeneity of variance was checked with assumption checks (Levene p =
0.569) and approximate normality of the residual state (ShapiroWilk p =
0.427). These differences are depicted in Fig 2 qualitatively where
situation awareness was lowest when the participants were in the Extreme
Hare environments and both the Ambiguous as well as the Extreme Stag
conditions did encourage high situation awareness. This trend indicates
that uncertainty of the environment (O) or high possible payoff (R)
enhances the monitoring and mental involvement of the participants,
which is in line with the theoretical perspective of viewed situation
awareness elevation with an increase in task demands or potential
payoffs (Fang et al., 2025).

\begin{Shaded}
\begin{Highlighting}[]
\CommentTok{\# ==============================}
\CommentTok{\# Table 2: Descriptive Statistics for Main Variables}
\CommentTok{\# ==============================}

\FunctionTok{library}\NormalTok{(dplyr)}
\FunctionTok{library}\NormalTok{(knitr)}
\FunctionTok{library}\NormalTok{(kableExtra)}

\CommentTok{\# 1. Workload (NASA{-}TLX total)}
\NormalTok{workload\_summary }\OtherTok{\textless{}{-}}\NormalTok{ workload }\SpecialCharTok{\%\textgreater{}\%}
  \FunctionTok{group\_by}\NormalTok{(condition) }\SpecialCharTok{\%\textgreater{}\%}
  \FunctionTok{summarise}\NormalTok{(}
    \AttributeTok{n =} \FunctionTok{n}\NormalTok{(),}
    \AttributeTok{Workload\_Mean =} \FunctionTok{round}\NormalTok{(}\FunctionTok{mean}\NormalTok{(total\_score\_sum\_6, }\AttributeTok{na.rm =} \ConstantTok{TRUE}\NormalTok{), }\DecValTok{1}\NormalTok{),}
    \AttributeTok{Workload\_SD   =} \FunctionTok{round}\NormalTok{(}\FunctionTok{sd}\NormalTok{(total\_score\_sum\_6, }\AttributeTok{na.rm =} \ConstantTok{TRUE}\NormalTok{), }\DecValTok{1}\NormalTok{),}
    \AttributeTok{.groups =} \StringTok{\textquotesingle{}drop\textquotesingle{}}
\NormalTok{  )}

\CommentTok{\# 2. NASA{-}TLX Subscales}
\NormalTok{subscale\_summary }\OtherTok{\textless{}{-}}\NormalTok{ workload }\SpecialCharTok{\%\textgreater{}\%}
  \FunctionTok{group\_by}\NormalTok{(condition) }\SpecialCharTok{\%\textgreater{}\%}
  \FunctionTok{summarise}\NormalTok{(}
    \AttributeTok{Mental\_Demand\_M =} \FunctionTok{round}\NormalTok{(}\FunctionTok{mean}\NormalTok{(mental\_demand, }\AttributeTok{na.rm =} \ConstantTok{TRUE}\NormalTok{),}\DecValTok{1}\NormalTok{),}
    \AttributeTok{Mental\_Demand\_SD =} \FunctionTok{round}\NormalTok{(}\FunctionTok{sd}\NormalTok{(mental\_demand, }\AttributeTok{na.rm =} \ConstantTok{TRUE}\NormalTok{),}\DecValTok{1}\NormalTok{),}
    \AttributeTok{Physical\_Demand\_M =} \FunctionTok{round}\NormalTok{(}\FunctionTok{mean}\NormalTok{(physical\_demand, }\AttributeTok{na.rm =} \ConstantTok{TRUE}\NormalTok{),}\DecValTok{1}\NormalTok{),}
    \AttributeTok{Physical\_Demand\_SD =} \FunctionTok{round}\NormalTok{(}\FunctionTok{sd}\NormalTok{(physical\_demand, }\AttributeTok{na.rm =} \ConstantTok{TRUE}\NormalTok{),}\DecValTok{1}\NormalTok{),}
    \AttributeTok{Effort\_M =} \FunctionTok{round}\NormalTok{(}\FunctionTok{mean}\NormalTok{(effort, }\AttributeTok{na.rm =} \ConstantTok{TRUE}\NormalTok{),}\DecValTok{1}\NormalTok{),}
    \AttributeTok{Effort\_SD =} \FunctionTok{round}\NormalTok{(}\FunctionTok{sd}\NormalTok{(effort, }\AttributeTok{na.rm =} \ConstantTok{TRUE}\NormalTok{),}\DecValTok{1}\NormalTok{),}
    \AttributeTok{Performance\_M =} \FunctionTok{round}\NormalTok{(}\FunctionTok{mean}\NormalTok{(performance, }\AttributeTok{na.rm =} \ConstantTok{TRUE}\NormalTok{),}\DecValTok{1}\NormalTok{),}
    \AttributeTok{Performance\_SD =} \FunctionTok{round}\NormalTok{(}\FunctionTok{sd}\NormalTok{(performance, }\AttributeTok{na.rm =} \ConstantTok{TRUE}\NormalTok{),}\DecValTok{1}\NormalTok{),}
    \AttributeTok{Frustration\_M =} \FunctionTok{round}\NormalTok{(}\FunctionTok{mean}\NormalTok{(frustration, }\AttributeTok{na.rm =} \ConstantTok{TRUE}\NormalTok{),}\DecValTok{1}\NormalTok{),}
    \AttributeTok{Frustration\_SD =} \FunctionTok{round}\NormalTok{(}\FunctionTok{sd}\NormalTok{(frustration, }\AttributeTok{na.rm =} \ConstantTok{TRUE}\NormalTok{),}\DecValTok{1}\NormalTok{),}
    \AttributeTok{.groups =} \StringTok{\textquotesingle{}drop\textquotesingle{}}
\NormalTok{  )}

\CommentTok{\# 3. Situation Awareness (SA total)}
\NormalTok{sa\_summary }\OtherTok{\textless{}{-}}\NormalTok{ sa }\SpecialCharTok{\%\textgreater{}\%}
  \FunctionTok{group\_by}\NormalTok{(demand\_condition) }\SpecialCharTok{\%\textgreater{}\%}
  \FunctionTok{summarise}\NormalTok{(}
    \AttributeTok{SA\_Mean =} \FunctionTok{round}\NormalTok{(}\FunctionTok{mean}\NormalTok{(sa\_total, }\AttributeTok{na.rm =} \ConstantTok{TRUE}\NormalTok{), }\DecValTok{1}\NormalTok{),}
    \AttributeTok{SA\_SD =} \FunctionTok{round}\NormalTok{(}\FunctionTok{sd}\NormalTok{(sa\_total, }\AttributeTok{na.rm =} \ConstantTok{TRUE}\NormalTok{), }\DecValTok{1}\NormalTok{),}
    \AttributeTok{.groups =} \StringTok{\textquotesingle{}drop\textquotesingle{}}
\NormalTok{  ) }\SpecialCharTok{\%\textgreater{}\%}
  \FunctionTok{rename}\NormalTok{(}\AttributeTok{condition =}\NormalTok{ demand\_condition)}

\CommentTok{\# 4. Performance (FullData)}
\NormalTok{performance\_summary }\OtherTok{\textless{}{-}}\NormalTok{ fulldata }\SpecialCharTok{\%\textgreater{}\%}
  \FunctionTok{group\_by}\NormalTok{(condition) }\SpecialCharTok{\%\textgreater{}\%}
  \FunctionTok{summarise}\NormalTok{(}
    \AttributeTok{Performance\_Mean =} \FunctionTok{round}\NormalTok{(}\FunctionTok{mean}\NormalTok{(score, }\AttributeTok{na.rm =} \ConstantTok{TRUE}\NormalTok{),}\DecValTok{1}\NormalTok{),}
    \AttributeTok{Performance\_SD   =} \FunctionTok{round}\NormalTok{(}\FunctionTok{sd}\NormalTok{(score, }\AttributeTok{na.rm =} \ConstantTok{TRUE}\NormalTok{),}\DecValTok{1}\NormalTok{),}
    \AttributeTok{.groups =} \StringTok{\textquotesingle{}drop\textquotesingle{}}
\NormalTok{  )}

\CommentTok{\# ==============================}
\CommentTok{\# Merge all summaries into one table}
\CommentTok{\# ==============================}
\NormalTok{table2 }\OtherTok{\textless{}{-}}\NormalTok{ workload\_summary }\SpecialCharTok{\%\textgreater{}\%}
  \FunctionTok{left\_join}\NormalTok{(subscale\_summary, }\AttributeTok{by =} \StringTok{"condition"}\NormalTok{) }\SpecialCharTok{\%\textgreater{}\%}
  \FunctionTok{left\_join}\NormalTok{(sa\_summary, }\AttributeTok{by =} \StringTok{"condition"}\NormalTok{) }\SpecialCharTok{\%\textgreater{}\%}
  \FunctionTok{left\_join}\NormalTok{(performance\_summary, }\AttributeTok{by =} \StringTok{"condition"}\NormalTok{)}

\CommentTok{\# Print nicely formatted table}
\NormalTok{table2 }\SpecialCharTok{\%\textgreater{}\%}
  \FunctionTok{kable}\NormalTok{(}
    \AttributeTok{caption =} \StringTok{"Table 2: Descriptive Statistics for Main Variables by Condition"}\NormalTok{,}
    \AttributeTok{booktabs =} \ConstantTok{TRUE}\NormalTok{,}
    \AttributeTok{digits =} \DecValTok{1}
\NormalTok{  ) }\SpecialCharTok{\%\textgreater{}\%}
  \FunctionTok{kable\_styling}\NormalTok{(}\AttributeTok{full\_width =} \ConstantTok{FALSE}\NormalTok{, }\AttributeTok{position =} \StringTok{"center"}\NormalTok{)}
\end{Highlighting}
\end{Shaded}

\begin{longtable}[t]{lrrrrrrrrrrrrrrrrr}
\caption{\label{tab:table2-descriptives}Table 2: Descriptive Statistics for Main Variables by Condition}\\
\toprule
condition & n & Workload\_Mean & Workload\_SD & Mental\_Demand\_M & Mental\_Demand\_SD & Physical\_Demand\_M & Physical\_Demand\_SD & Effort\_M & Effort\_SD & Performance\_M & Performance\_SD.x & Frustration\_M & Frustration\_SD & SA\_Mean & SA\_SD & Performance\_Mean & Performance\_SD.y\\
\midrule
E & 10 & 46.2 & 15.1 & 57 & 27.8 & 28.5 & 17.5 & 59 & 17.4 & 42.5 & 26.4 & 33.0 & 21.0 & 10.2 & 4.2 & 691.2 & 88.8\\
O & 10 & 38.1 & 14.8 & 54 & 19.3 & 16.5 & 14.9 & 52 & 21.8 & 25.5 & 11.4 & 32.5 & 21.0 & 15.1 & 2.9 & 751.5 & 95.7\\
R & 10 & 34.2 & 11.0 & 48 & 24.5 & 13.0 & 9.2 & 34 & 18.5 & 32.0 & 19.0 & 25.5 & 17.4 & 14.6 & 3.9 & 795.6 & 36.0\\
\bottomrule
\end{longtable}

\begin{Shaded}
\begin{Highlighting}[]
\FunctionTok{library}\NormalTok{(ggplot2)}
\FunctionTok{library}\NormalTok{(dplyr)}

\CommentTok{\# Prepare data for plotting}
\NormalTok{sa\_plot\_data }\OtherTok{\textless{}{-}}\NormalTok{ sa }\SpecialCharTok{\%\textgreater{}\%}
  \FunctionTok{select}\NormalTok{(participant, demand\_condition, sa\_total) }\SpecialCharTok{\%\textgreater{}\%}
  \FunctionTok{arrange}\NormalTok{(demand\_condition, participant)}

\CommentTok{\# Summary statistics for plotting}
\NormalTok{sa\_summary }\OtherTok{\textless{}{-}}\NormalTok{ sa\_plot\_data }\SpecialCharTok{\%\textgreater{}\%}
  \FunctionTok{group\_by}\NormalTok{(demand\_condition) }\SpecialCharTok{\%\textgreater{}\%}
  \FunctionTok{summarise}\NormalTok{(}
    \AttributeTok{n =} \FunctionTok{n}\NormalTok{(),}
    \AttributeTok{mean\_sa =} \FunctionTok{mean}\NormalTok{(sa\_total, }\AttributeTok{na.rm =} \ConstantTok{TRUE}\NormalTok{),}
    \AttributeTok{sd\_sa =} \FunctionTok{sd}\NormalTok{(sa\_total, }\AttributeTok{na.rm =} \ConstantTok{TRUE}\NormalTok{),}
    \AttributeTok{se\_sa =}\NormalTok{ sd\_sa }\SpecialCharTok{/} \FunctionTok{sqrt}\NormalTok{(n),}
    \AttributeTok{.groups =} \StringTok{\textquotesingle{}drop\textquotesingle{}}
\NormalTok{  )}

\CommentTok{\# Print summary table}
\NormalTok{sa\_summary}
\end{Highlighting}
\end{Shaded}

\begin{verbatim}
## # A tibble: 3 x 5
##   demand_condition     n mean_sa sd_sa se_sa
##   <fct>            <int>   <dbl> <dbl> <dbl>
## 1 E                   10    10.2  4.24 1.34 
## 2 O                   10    15.1  2.88 0.912
## 3 R                   10    14.6  3.92 1.24
\end{verbatim}

\begin{Shaded}
\begin{Highlighting}[]
\CommentTok{\# Create boxplot with individual data points}
\FunctionTok{ggplot}\NormalTok{(sa\_plot\_data, }\FunctionTok{aes}\NormalTok{(}\AttributeTok{x =}\NormalTok{ demand\_condition, }\AttributeTok{y =}\NormalTok{ sa\_total, }\AttributeTok{fill =}\NormalTok{ demand\_condition)) }\SpecialCharTok{+}
  \FunctionTok{geom\_boxplot}\NormalTok{(}\AttributeTok{alpha =} \FloatTok{0.7}\NormalTok{, }\AttributeTok{outlier.shape =} \ConstantTok{NA}\NormalTok{) }\SpecialCharTok{+}
  \FunctionTok{geom\_jitter}\NormalTok{(}\AttributeTok{width =} \FloatTok{0.2}\NormalTok{, }\AttributeTok{alpha =} \FloatTok{0.4}\NormalTok{, }\AttributeTok{color =} \StringTok{"black"}\NormalTok{) }\SpecialCharTok{+}
  \FunctionTok{theme\_minimal}\NormalTok{(}\AttributeTok{base\_size =} \DecValTok{12}\NormalTok{) }\SpecialCharTok{+}
  \FunctionTok{labs}\NormalTok{(}
    \AttributeTok{title =} \StringTok{"Figure 2: Situation Awareness Total Score by Condition"}\NormalTok{,}
    \AttributeTok{x =} \StringTok{"Demand Condition"}\NormalTok{,}
    \AttributeTok{y =} \StringTok{"SA Total Score"}
\NormalTok{  ) }\SpecialCharTok{+}
  \FunctionTok{theme}\NormalTok{(}\AttributeTok{legend.position =} \StringTok{"none"}\NormalTok{,}
        \AttributeTok{axis.text.x =} \FunctionTok{element\_text}\NormalTok{(}\AttributeTok{angle =} \DecValTok{45}\NormalTok{, }\AttributeTok{hjust =} \DecValTok{1}\NormalTok{)) }\SpecialCharTok{+}
  \FunctionTok{scale\_fill\_brewer}\NormalTok{(}\AttributeTok{palette =} \StringTok{"Set2"}\NormalTok{)}
\end{Highlighting}
\end{Shaded}

\includegraphics{YY56_files/figure-latex/figure2-sa-total-1.pdf} \#\#\#
Performance (Score) The scores of performance varied significantly
across conditions, F (2,27) =4.492, p=0.021. A Mean difference of 104.4,
95\% CI {[}17.69, 191.11{]} at p = 0.016 suggested that the scores in
the Extreme\_ Stag (R) group were significantly higher than in the
Extreme\_Hare (E) group. The results did not provide significant
differences between Ambiguous (O) and E (diff = 60.3, p = 0.215) and R
vs O (diff = 44.1, p = 0.429) indicating that the Stag team had a strong
performance advantage only in the circumstances when the environment was
conducive to the game. There was no issue with assumption checks,
(Levene p = 0.269; ShapiroWilk p = 0.091). The aggregate and individual
performance data are summarised in figure 3. Mean scores grew between E
(M = 691, SD = 88.8) and O (M = 752, SD = 95.7), and were maximum in R
(M = 796, SD = 36.0). The means were also pretty close to the median
values and the ranges were smaller in R, which implied greater
consistency in performance under the Stag-favouring condition. The
figure shows an increase in rewards in participants with clear
environmental structure and high-payoff option which are significant in
the role of environmental cues on the effectiveness of their decisions
and the choice of strategy (Bollen et al., 2025).

\begin{Shaded}
\begin{Highlighting}[]
\FunctionTok{library}\NormalTok{(ggplot2)}
\FunctionTok{library}\NormalTok{(dplyr)}

\CommentTok{\# Prepare data for plotting}
\NormalTok{performance\_plot\_data }\OtherTok{\textless{}{-}}\NormalTok{ fulldata }\SpecialCharTok{\%\textgreater{}\%}
  \FunctionTok{select}\NormalTok{(participant, condition, score) }\SpecialCharTok{\%\textgreater{}\%}
  \FunctionTok{arrange}\NormalTok{(condition, participant)}

\CommentTok{\# Summary statistics for reference}
\NormalTok{performance\_summary }\OtherTok{\textless{}{-}}\NormalTok{ performance\_plot\_data }\SpecialCharTok{\%\textgreater{}\%}
  \FunctionTok{group\_by}\NormalTok{(condition) }\SpecialCharTok{\%\textgreater{}\%}
  \FunctionTok{summarise}\NormalTok{(}
    \AttributeTok{n =} \FunctionTok{n}\NormalTok{(),}
    \AttributeTok{mean\_score =} \FunctionTok{mean}\NormalTok{(score, }\AttributeTok{na.rm =} \ConstantTok{TRUE}\NormalTok{),}
    \AttributeTok{sd\_score =} \FunctionTok{sd}\NormalTok{(score, }\AttributeTok{na.rm =} \ConstantTok{TRUE}\NormalTok{),}
    \AttributeTok{se\_score =}\NormalTok{ sd\_score }\SpecialCharTok{/} \FunctionTok{sqrt}\NormalTok{(n),}
    \AttributeTok{.groups =} \StringTok{\textquotesingle{}drop\textquotesingle{}}
\NormalTok{  )}

\CommentTok{\# Print summary table}
\NormalTok{performance\_summary}
\end{Highlighting}
\end{Shaded}

\begin{verbatim}
## # A tibble: 3 x 5
##   condition     n mean_score sd_score se_score
##   <fct>     <int>      <dbl>    <dbl>    <dbl>
## 1 E            10       691.     88.8     28.1
## 2 O            10       752.     95.7     30.3
## 3 R            10       796.     36.0     11.4
\end{verbatim}

\begin{Shaded}
\begin{Highlighting}[]
\CommentTok{\# Create violin + boxplot with individual data points}
\FunctionTok{ggplot}\NormalTok{(performance\_plot\_data, }\FunctionTok{aes}\NormalTok{(}\AttributeTok{x =}\NormalTok{ condition, }\AttributeTok{y =}\NormalTok{ score, }\AttributeTok{fill =}\NormalTok{ condition)) }\SpecialCharTok{+}
  \FunctionTok{geom\_violin}\NormalTok{(}\AttributeTok{alpha =} \FloatTok{0.6}\NormalTok{, }\AttributeTok{trim =} \ConstantTok{TRUE}\NormalTok{) }\SpecialCharTok{+}
  \FunctionTok{geom\_boxplot}\NormalTok{(}\AttributeTok{width =} \FloatTok{0.2}\NormalTok{, }\AttributeTok{alpha =} \FloatTok{0.8}\NormalTok{, }\AttributeTok{outlier.shape =} \ConstantTok{NA}\NormalTok{) }\SpecialCharTok{+}
  \FunctionTok{geom\_jitter}\NormalTok{(}\AttributeTok{width =} \FloatTok{0.15}\NormalTok{, }\AttributeTok{alpha =} \FloatTok{0.4}\NormalTok{, }\AttributeTok{color =} \StringTok{"black"}\NormalTok{) }\SpecialCharTok{+}
  \FunctionTok{theme\_minimal}\NormalTok{(}\AttributeTok{base\_size =} \DecValTok{12}\NormalTok{) }\SpecialCharTok{+}
  \FunctionTok{labs}\NormalTok{(}
    \AttributeTok{title =} \StringTok{"Figure 3: Performance Scores by Condition"}\NormalTok{,}
    \AttributeTok{x =} \StringTok{"Condition"}\NormalTok{,}
    \AttributeTok{y =} \StringTok{"Performance Score"}
\NormalTok{  ) }\SpecialCharTok{+}
  \FunctionTok{theme}\NormalTok{(}\AttributeTok{legend.position =} \StringTok{"none"}\NormalTok{,}
        \AttributeTok{axis.text.x =} \FunctionTok{element\_text}\NormalTok{(}\AttributeTok{angle =} \DecValTok{45}\NormalTok{, }\AttributeTok{hjust =} \DecValTok{1}\NormalTok{)) }\SpecialCharTok{+}
  \FunctionTok{scale\_fill\_brewer}\NormalTok{(}\AttributeTok{palette =} \StringTok{"Pastel1"}\NormalTok{)}
\end{Highlighting}
\end{Shaded}

\includegraphics{YY56_files/figure-latex/figure3-performance-1.pdf}

\subsubsection{NASA-TLX subscales}\label{nasa-tlx-subscales}

Individual subscale analysis of NASA-TLX showed that there was
significant difference in reported effort concerning conditions, F(2,27)
= 4.451, p = 0.021, with the participants in Extreme\_Hare (E), and
Ambiguous (O) conditions reporting higher effort levels (M = 59 and 52,
respectively) as compared to the Extreme\_Stag (R) condition (M = 34).
The physical demand was also significant, F(2,27) = 3.235, p = 0.055, E
had the largest mean (28.5) and R the lowest (13). Mental demand,
temporal demand, performance and frustration did not differ
significantly, as did other subscales. These patterns are shown in
figure 4 with higher effort and consequently, lower physical demand
being used in higher uncertainty or risk conditions and lower subjective
workload in the Stag-favouring environment. The standard errors and the
means show that there are always differences, especially the lower mean
and the lower variability of effort as shown by R (Bollen et al., 2025).
This trend indicates that the Stag setting was less challenging to the
participants, probably because of better task organization and more
predictable results, and ambiguous and Hare settings had to be attended
to and closely monitored due to their ambiguity, which added perceived
effort and physical load.

\begin{Shaded}
\begin{Highlighting}[]
\FunctionTok{library}\NormalTok{(ggplot2)}
\FunctionTok{library}\NormalTok{(dplyr)}
\FunctionTok{library}\NormalTok{(tidyr)}

\CommentTok{\# Reshape NASA{-}TLX subscale data for plotting}
\NormalTok{workload\_melt }\OtherTok{\textless{}{-}}\NormalTok{ workload }\SpecialCharTok{\%\textgreater{}\%}
  \FunctionTok{select}\NormalTok{(participant, condition, mental\_demand, physical\_demand, }
\NormalTok{         temporal\_demand, performance, effort, frustration) }\SpecialCharTok{\%\textgreater{}\%}
  \FunctionTok{pivot\_longer}\NormalTok{(}
    \AttributeTok{cols =}\NormalTok{ mental\_demand}\SpecialCharTok{:}\NormalTok{frustration,}
    \AttributeTok{names\_to =} \StringTok{"subscale"}\NormalTok{,}
    \AttributeTok{values\_to =} \StringTok{"score"}
\NormalTok{  )}

\CommentTok{\# Calculate mean and SE for each subscale x condition}
\NormalTok{subscale\_summary }\OtherTok{\textless{}{-}}\NormalTok{ workload\_melt }\SpecialCharTok{\%\textgreater{}\%}
  \FunctionTok{group\_by}\NormalTok{(condition, subscale) }\SpecialCharTok{\%\textgreater{}\%}
  \FunctionTok{summarise}\NormalTok{(}
    \AttributeTok{n =} \FunctionTok{n}\NormalTok{(),}
    \AttributeTok{mean\_score =} \FunctionTok{mean}\NormalTok{(score, }\AttributeTok{na.rm =} \ConstantTok{TRUE}\NormalTok{),}
    \AttributeTok{sd\_score =} \FunctionTok{sd}\NormalTok{(score, }\AttributeTok{na.rm =} \ConstantTok{TRUE}\NormalTok{),}
    \AttributeTok{se\_score =}\NormalTok{ sd\_score }\SpecialCharTok{/} \FunctionTok{sqrt}\NormalTok{(n),}
    \AttributeTok{.groups =} \StringTok{\textquotesingle{}drop\textquotesingle{}}
\NormalTok{  )}

\CommentTok{\# Print summary table}
\NormalTok{subscale\_summary}
\end{Highlighting}
\end{Shaded}

\begin{verbatim}
## # A tibble: 18 x 6
##    condition subscale            n mean_score sd_score se_score
##    <fct>     <chr>           <int>      <dbl>    <dbl>    <dbl>
##  1 E         effort             10       59      17.4      5.52
##  2 E         frustration        10       33      21.0      6.63
##  3 E         mental_demand      10       57      27.8      8.79
##  4 E         performance        10       42.5    26.4      8.34
##  5 E         physical_demand    10       28.5    17.5      5.53
##  6 E         temporal_demand    10       57      19.9      6.29
##  7 O         effort             10       52      21.8      6.88
##  8 O         frustration        10       32.5    21.0      6.64
##  9 O         mental_demand      10       54      19.3      6.09
## 10 O         performance        10       25.5    11.4      3.61
## 11 O         physical_demand    10       16.5    14.9      4.72
## 12 O         temporal_demand    10       48      27.1      8.57
## 13 R         effort             10       34      18.5      5.86
## 14 R         frustration        10       25.5    17.4      5.5 
## 15 R         mental_demand      10       48      24.5      7.75
## 16 R         performance        10       32      19.0      6.02
## 17 R         physical_demand    10       13       9.19     2.91
## 18 R         temporal_demand    10       52.6    18.4      5.83
\end{verbatim}

\begin{Shaded}
\begin{Highlighting}[]
\CommentTok{\# Create grouped barplot with error bars}
\FunctionTok{ggplot}\NormalTok{(subscale\_summary, }\FunctionTok{aes}\NormalTok{(}\AttributeTok{x =}\NormalTok{ subscale, }\AttributeTok{y =}\NormalTok{ mean\_score, }\AttributeTok{fill =}\NormalTok{ condition)) }\SpecialCharTok{+}
  \FunctionTok{geom\_bar}\NormalTok{(}\AttributeTok{stat =} \StringTok{"identity"}\NormalTok{, }\AttributeTok{position =} \FunctionTok{position\_dodge}\NormalTok{(}\AttributeTok{width =} \FloatTok{0.8}\NormalTok{), }\AttributeTok{alpha =} \FloatTok{0.8}\NormalTok{) }\SpecialCharTok{+}
  \FunctionTok{geom\_errorbar}\NormalTok{(}\FunctionTok{aes}\NormalTok{(}\AttributeTok{ymin =}\NormalTok{ mean\_score }\SpecialCharTok{{-}}\NormalTok{ se\_score, }\AttributeTok{ymax =}\NormalTok{ mean\_score }\SpecialCharTok{+}\NormalTok{ se\_score),}
                \AttributeTok{width =} \FloatTok{0.2}\NormalTok{, }\AttributeTok{position =} \FunctionTok{position\_dodge}\NormalTok{(}\FloatTok{0.8}\NormalTok{)) }\SpecialCharTok{+}
  \FunctionTok{theme\_minimal}\NormalTok{(}\AttributeTok{base\_size =} \DecValTok{12}\NormalTok{) }\SpecialCharTok{+}
  \FunctionTok{labs}\NormalTok{(}
    \AttributeTok{title =} \StringTok{"Figure 4: Mean NASA{-}TLX Subscale Scores by Condition"}\NormalTok{,}
    \AttributeTok{x =} \StringTok{"NASA{-}TLX Subscale"}\NormalTok{,}
    \AttributeTok{y =} \StringTok{"Mean Score"}
\NormalTok{  ) }\SpecialCharTok{+}
  \FunctionTok{theme}\NormalTok{(}\AttributeTok{axis.text.x =} \FunctionTok{element\_text}\NormalTok{(}\AttributeTok{angle =} \DecValTok{45}\NormalTok{, }\AttributeTok{hjust =} \DecValTok{1}\NormalTok{)) }\SpecialCharTok{+}
  \FunctionTok{scale\_fill\_brewer}\NormalTok{(}\AttributeTok{palette =} \StringTok{"Pastel2"}\NormalTok{)}
\end{Highlighting}
\end{Shaded}

\includegraphics{YY56_files/figure-latex/figure4-nasatlx-1.pdf} \#\#\#
Correlation Pearson correlation studies found out that there were
various meaningful correlations among continuous variables. Choice
accuracy had a strong positive correlation with performance (score)
which = 0.640, p = 0.001, which means that participants who made more
accurate choices of the team earned higher rewards. Response time was
also found to have a negative correlation with the choice accuracy, r =
-0.446, p = 0.014, meaning that the quicker the choice was the more
accurate the choice was chosen. There was a moderate but negative
relationship between NASA-TLX total and SART total, r = - 0.375, p =
0.04 1, which suggested that the greater perceived workload was, the
lower was the situation awareness. Other correlations were not
significant or weak performance and NASA-TLX, r = -0.120; performance
and SART, r = 0.014; NASA-TLX and choice accuracy, r = -0.263; SART and
choice accuracy, r = 0.060; and response time and score, r = -0.156 and
SART, r = -0.108. These patterns have been visualised in figure 5, and
have the strongest associations between the performance, accuracy, and
response efficiency, and the negative relation between the workload and
situation awareness, which correlates to the cognitive load theory
(Soetewey, 2020).

\begin{Shaded}
\begin{Highlighting}[]
\FunctionTok{library}\NormalTok{(ggplot2)}
\FunctionTok{library}\NormalTok{(reshape2)}
\FunctionTok{library}\NormalTok{(dplyr)}
\FunctionTok{library}\NormalTok{(psych)}

\CommentTok{\# Select relevant continuous variables for correlation}
\NormalTok{cor\_data }\OtherTok{\textless{}{-}}\NormalTok{ fulldata }\SpecialCharTok{\%\textgreater{}\%}
  \FunctionTok{select}\NormalTok{(score, nasa\_tlx, sart, choice\_accuracy, time\_respond\_env)}

\CommentTok{\# Compute correlation matrix}
\NormalTok{cor\_matrix }\OtherTok{\textless{}{-}} \FunctionTok{cor}\NormalTok{(cor\_data, }\AttributeTok{use =} \StringTok{"pairwise.complete.obs"}\NormalTok{)}

\CommentTok{\# Melt correlation matrix for plotting}
\NormalTok{cor\_melt }\OtherTok{\textless{}{-}} \FunctionTok{melt}\NormalTok{(cor\_matrix)}

\CommentTok{\# Print correlation values}
\NormalTok{cor\_melt}
\end{Highlighting}
\end{Shaded}

\begin{verbatim}
##                Var1             Var2       value
## 1             score            score  1.00000000
## 2          nasa_tlx            score -0.12034368
## 3              sart            score  0.01374808
## 4   choice_accuracy            score  0.64049580
## 5  time_respond_env            score -0.15601108
## 6             score         nasa_tlx -0.12034368
## 7          nasa_tlx         nasa_tlx  1.00000000
## 8              sart         nasa_tlx -0.37455400
## 9   choice_accuracy         nasa_tlx -0.26335153
## 10 time_respond_env         nasa_tlx  0.28790328
## 11            score             sart  0.01374808
## 12         nasa_tlx             sart -0.37455400
## 13             sart             sart  1.00000000
## 14  choice_accuracy             sart  0.06041736
## 15 time_respond_env             sart -0.10828335
## 16            score  choice_accuracy  0.64049580
## 17         nasa_tlx  choice_accuracy -0.26335153
## 18             sart  choice_accuracy  0.06041736
## 19  choice_accuracy  choice_accuracy  1.00000000
## 20 time_respond_env  choice_accuracy -0.44590621
## 21            score time_respond_env -0.15601108
## 22         nasa_tlx time_respond_env  0.28790328
## 23             sart time_respond_env -0.10828335
## 24  choice_accuracy time_respond_env -0.44590621
## 25 time_respond_env time_respond_env  1.00000000
\end{verbatim}

\begin{Shaded}
\begin{Highlighting}[]
\CommentTok{\# Create correlation heatmap}
\FunctionTok{ggplot}\NormalTok{(cor\_melt, }\FunctionTok{aes}\NormalTok{(}\AttributeTok{x =}\NormalTok{ Var1, }\AttributeTok{y =}\NormalTok{ Var2, }\AttributeTok{fill =}\NormalTok{ value)) }\SpecialCharTok{+}
  \FunctionTok{geom\_tile}\NormalTok{(}\AttributeTok{color =} \StringTok{"white"}\NormalTok{) }\SpecialCharTok{+}
  \FunctionTok{geom\_text}\NormalTok{(}\FunctionTok{aes}\NormalTok{(}\AttributeTok{label =} \FunctionTok{round}\NormalTok{(value, }\DecValTok{2}\NormalTok{)), }\AttributeTok{color =} \StringTok{"black"}\NormalTok{, }\AttributeTok{size =} \DecValTok{3}\NormalTok{) }\SpecialCharTok{+}
  \FunctionTok{scale\_fill\_gradient2}\NormalTok{(}\AttributeTok{low =} \StringTok{"blue"}\NormalTok{, }\AttributeTok{mid =} \StringTok{"white"}\NormalTok{, }\AttributeTok{high =} \StringTok{"red"}\NormalTok{, }
                       \AttributeTok{midpoint =} \DecValTok{0}\NormalTok{, }\AttributeTok{limit =} \FunctionTok{c}\NormalTok{(}\SpecialCharTok{{-}}\DecValTok{1}\NormalTok{, }\DecValTok{1}\NormalTok{), }\AttributeTok{name =} \StringTok{"Correlation"}\NormalTok{) }\SpecialCharTok{+}
  \FunctionTok{theme\_minimal}\NormalTok{(}\AttributeTok{base\_size =} \DecValTok{12}\NormalTok{) }\SpecialCharTok{+}
  \FunctionTok{labs}\NormalTok{(}
    \AttributeTok{title =} \StringTok{"Figure 5: Correlation Matrix of Performance Measures"}\NormalTok{,}
    \AttributeTok{x =} \StringTok{""}\NormalTok{,}
    \AttributeTok{y =} \StringTok{""}
\NormalTok{  ) }\SpecialCharTok{+}
  \FunctionTok{theme}\NormalTok{(}\AttributeTok{axis.text.x =} \FunctionTok{element\_text}\NormalTok{(}\AttributeTok{angle =} \DecValTok{45}\NormalTok{, }\AttributeTok{hjust =} \DecValTok{1}\NormalTok{),}
        \AttributeTok{axis.text.y =} \FunctionTok{element\_text}\NormalTok{(}\AttributeTok{hjust =} \DecValTok{1}\NormalTok{))}
\end{Highlighting}
\end{Shaded}

\includegraphics{YY56_files/figure-latex/figure5-correlation-1.pdf}

\subsubsection{Reliability}\label{reliability}

This study showed an acceptable level of internal consistency of
NASA-TLX subscales in its form, with a Cronbach alpha of 0.78 and a 95\%
interval of 0.64 to 0.88. This reliability is a good indication of the
internal consistency of the six NASA-TLX subscales and hence they can be
used as a valid measure of subjective workload among the sample. In line
with accepted principles, an alpha of over 0.70 implies that the
instrument is able to measure a single construct, which is perceived
workload. These results are similar to the previous literature
indicating that NASA-TLX is a strong and extensively utilized workload
measurement tool in the human-machine interaction research (Kotian,
Varghese and Motappa, 2022).

\section{Discussion}\label{discussion}

\subsection{Summary and interpretation of key
findings}\label{summary-and-interpretation-of-key-findings}

In this study, the question presented was whether or not human
supervisors could be effective in choosing between specialist and
adaptive multi-agent reinforcement learning (MARL) teams in environments
of varying risk profile. According to the results, there is a definite
performance advantage in the environments where high-reward strategies
are strongly provided. To be more specific, the greatest mean
performance score was registered in Extreme\_Stag (R) condition (M =
795.6) which was significantly lower than the Extreme\_Hare (E)
condition with the mean difference of 104.4 (p = 0.0159). It implies
that in the case when an environment intensively rewards a high-reward
alternative, supervisors can better identify and take advantage of the
underlying task structure or can harness the inherent efficiency of the
Stag-specialist team (Yang, He and Huang, 2024). The situation awareness
was varying systematically in environmental situations. There were
significant variations in SART total scores between the Extreme\_Hare
(E) and the Ambiguous (O) conditions and between the Extreme\_Hare (E)
and the Extreme Stag (R) conditions. The implication of these findings
is that there is an increment in supervisory attention and awareness in
uncertain conditions, or where there is saliency of potential rewards.
On the other hand, the Extreme\_Hare state was linked to less situation
awareness indicating that conditions with lesser incentives or cues can
diminish engagements and monitoring. The workload as a result of the
NASA-TLX total scores showed no significant difference between the
conditions (F(2,27) = 1.976, p = 0.158; BF = 0.733) which developed
moderate Bayesian support of the null hypothesis. Nonetheless, the
effort subscale did show significant differences (p = 0.021), and the
effect size between workload components seemed to be moderate-large when
the workload is needed to be considered constant between the
Extreme\_Stag and Extreme\_Hare (Cohen's d = 0.91), which indicates that
components of workload can change despite the apparent stability of the
workload (Griffiths et al., 2020). Lastly, performance was strongly
related to behavioural efficiency. There is a positive correlation
between performance and accuracy of choice (r = 0.64, p = 0.001; r =
0.446, p = 0.014) and speed of responding which attests to the
established relations between accuracy, efficiency and outcome (Sanchez
and De La Torre, 1996). Although the awareness level was higher in
unclear trials, supervisors did not always favor the adaptive DisAdv
team, which is also consistent with the fact that humans choose not to
trust adaptive systems without clear guidelines of their usefulness
(Vaccaro et al., 2021).

\subsection{Relation to existing
literature}\label{relation-to-existing-literature}

The performance benefits of the clear Stag settings (R) are consistent
with meta-analytic findings on the reduction in the cognitive load and
the improvement of decision outcomes when using unambiguous cues
(Onnasch et al., 2014). Respondents were more likely to respond
accurately and quickly in these settings, which is in line with the
observation that when the level of ambiguity reduces, more tasks are
performed. The high state of situational awareness can be found in the
ambiguous situations which are observed, which is in accordance with the
theoretical framework of Endsley (1995) that argues that uncertainty
requires more situational monitoring and information processing.
However, unlike the expectation that ambiguous conditions should cause
the participants to prefer adaptive agents, DisAdv were selected to a
lesser degree. This finding resembles the previous studies by Vaccaro et
al.~(2021) and Schmutz et al.~(2023), which proved that trust and
adaptive systems adoption are conditional on transparency, reason
clarity, and adaptive behaviours communication. The workload-situation
awareness negative correlation (r = -0.375, p = 0.041) confirms the
popular belief that high cognitive or perceived work load would
deteriorate situation awareness and quality of decisions (Hart and
Staveland, 1988). It means that the attentional resources of supervisors
can be overloaded in situations where more effort is required, which
might make them unable to track the states of systems effectively.

\subsection{Practical implications}\label{practical-implications}

A number of practical implications can be drawn out of these findings.
To start with, there is the requirement of transparency and adaptive
agent training. Clear signs of its detection behaviour and success rates
of the DisAdv agent can also be offered as a way of influencing
supervisors to use it more suitably in a grey setting (Sanchez and De La
Torre, 1996). Second, interface design must be focused on improving the
situation awareness and lessening the effort (Suriya and Nivetha, 2023).
As an example, the introduction of visuals like the number of berries
can be used as cues that can enhance the perception and precision of
environmental indicators that are not explicit. Third, work load sub
components should be evaluated in an experimental and operational way.
The effect of meaning differences in effort or physical demand that
affect behaviour and quality of decision making may be masked by
composite TLX scores.

\subsection{Limitations}\label{limitations}

There are a number of limitations that need to be mentioned. The sample
size of N = 30 is small, and this can also limit the statistical power
of detecting smaller effects (Mascha and Vetter, 2018). Some
hypothesized effects may also not have much evidence based on the
analysis of Bayesian. There was also a limit on the number of
mixed-effects modelling that could address repeated measures contained
in the participants because of aggregation of trial-level data. The
level of briefing on the capabilities of the agents especially to the
adaptive DisAdv team might have made a difference; there might have been
low transparency in the selection of adaptive agents. Lastly, the
laboratory setting and its unnatural generalizability of the laboratory
task, including the number of the berries and the probability of an
underlying illness, might restrict generalizability to real-world
complex supervisory contexts, and one should test ecological validity in
the future.

\subsection{Future work}\label{future-work}

The future study must offer concrete acts of transparency,
e.g.~rationale displays or post-selection feedback, to ascertain whether
supervisors are more effective in the selection of agents who are
adaptive in uncertain situations. Complete data at the trial-level would
allow hierarchical mixed-effects modelling to model both within-subject
and between-subject variance and investigate effects of time on learning
(Jost and Jansen, 2022). The assessment of interaction between agent
explanation, participant experience and team choice would be possible
with increasing sample size and different briefing or training
protocols. Lastly, incorporating eye-tracking or think-aloud procedures
may shed more light on the exact cues and approaches employed by
supervisors to make team selections, which can be used to inform the
interface design and adaptive agent development.

\#Conclusion This paper shows that human bosses work best in situations
characterized by clarity and Stag-favouring and situation awareness
changes with environmental uncertainty (The R Foundation, 2019). Even
though, ambiguous contexts evoked higher levels of reported SA, there
was no a priori reason why adaptive DisAdv team proved to be more
frequently chosen by supervisors. There was a positive relationship
between performance and choice accuracy, and the faster the response
time the accuracy. The immediate implication is that adaptive teams of
agents are not necessarily preferred; in order to enable human-AI
teaming, designers need to focus on transparency transparency, training,
and good interface indicators to enable supervisors to see the payoff of
adaptive agents. The following future study using trial-level data and
bigger sample sizes would enable more rigorous research of
supervisor-agent coordination strategies (Jost and Jansen, 2022 ).

\section{References}\label{references}

Bollen, Z., Wyss, A.M., Guizar Rosales, E., Le Houcq Corbi, Z. and
Knoch, D. (2025). The role of attitudes and attention in
pro-environmental decision-making: An eye-tracking study. Journal of
Environmental Psychology, 107, p.102776. doi:
\url{https://doi.org/10.1016/j.jenvp.2025.102776}. Collins, B.J.,
Munyon, T.P., Ashkanasy, N., Gallagher, E. and Gallagher, S.A. (2022).
Team Decision Making in Crisis: How Affect Asymmetry and Team Process
Norms Affect Decision-Making. Research on Emotion in Organizations.
{[}online{]} Available at:
\url{https://www.researchgate.net/publication/360082673_Team_Decision_Making_in_Crisis_How_Affect_Asymmetry_and_Team_Process_Norms_Affect_Decision-Making}.
Dung, T. and Duong, N.T. (2024). Investigating learning burnout and
academic performance among management students: a longitudinal study in
English courses. BMC psychology, 12(1). Doi:
\url{https://doi.org/10.1186/s40359-024-01725-6}.\\
Edwin De Jonge and Mark Van Der Loo (2013). Discussion Paper An
introduction to data cleaning with R. {[}online{]} Available at:
\url{https://cran.r-project.org/doc/contrib/de_Jonge+van_der_Loo-Introduction_to_data_cleaning_with_R.pdf}.
Emmanuelle Apartis (2013). Clinical neurophysiology in movement
disorders. Handbook of clinical neurology, {[}online{]} pp.87--92. doi:
\url{https://doi.org/10.1016/b978-0-444-52891-9.00008-7}. Endsley, M.R.
(1995). Toward a theory of situation awareness in dynamic systems. Human
Factors. {[}online{]} psycnet.apa.org. Available at:
\url{https://psycnet.apa.org/record/1995-42276-001}. Fang, R., Zhou, Q.,
Zhou, C., Yuan, S., Wang, K., Li, Q., Zhang, Y. and Li, J. (2025). A
Systematic Review of Empirical Studies on Situation Awareness:
Perspectives From the Interaction Among Humans, Machines, and the Task
Environment. PsyCh Journal, {[}online{]} 14(5), pp.718--733. doi:
\url{https://doi.org/10.1002/pchj.70027}. Griffiths, P., Saville, C.,
Ball, J., Jones, J., Pattison, N. and Monks, T. (2020). Nursing
workload, Nurse Staffing Methodologies \& Tools: a Systematic Scoping
Review \& Discussion. International Journal of Nursing Studies,
{[}online{]} 103(1), pp.1--11. doi:
\url{https://doi.org/10.1016/j.ijnurstu.2019.103487}.\\
Hart, S.G. and Staveland, L.E. (1988). Development of NASA-TLX (task
load index): Results of empirical and theoretical research. Advances in
Psychology, {[}online{]} 52, pp.139--183. doi:
\url{https://doi.org/10.1016/s0166-4115(08)62386-9}. Jost, L. and
Jansen, P. (2022). Using linear mixed models to analyze learning
processes within sessions improves detection of treatment effects: An
exemplary study of chronometric mental rotation. Methods in Psychology,
6, p.100092. doi: \url{https://doi.org/10.1016/j.metip.2022.100092}.
Juarros-Basterretxea, J., Aonso-Diego, G., Postigo, Á., Montes-Álvarez,
P., Menéndez-Aller, Á. and García-Cueto, E. (2024). Post-Hoc Tests in
One-Way ANOVA: The Case for Normal Distribution. Methodology,
{[}online{]} 20(2), pp.84--99. doi:
\url{https://doi.org/10.5964/meth.11721}. Koopmans, L., Coffeng, J.K.,
Bernaards, C.M., Boot, C.R., Hildebrandt, V.H., de Vet, H.C. and van der
Beek, A.J. (2014). Responsiveness of the individual work performance
questionnaire. BMC Public Health, 14(1). doi:
\url{https://doi.org/10.1186/1471-2458-14-513}. Kotian, H., Varghese,
A.L. and Motappa, R. (2022). An R Function for Cronbach's Alpha
Analysis: A Case-Based Approach. National journal of community medicine,
13(08), pp.571--575. doi:
\url{https://doi.org/10.55489/njcm.130820221149}. Lambdin, C. and
Shaffer, V.A. (2009). Are within-subjects designs transparent? Judgment
and Decision Making, 4(7), pp.544--566. doi:
\url{https://doi.org/10.1017/s1930297500001133}. Mascha, E.J. and
Vetter, T.R. (2018). Significance, Errors, Power, and Sample Size: The
Blocking and Tackling of Statistics. Anesthesia and Analgesia, 126(2),
pp.691--698. doi: \url{https://doi.org/10.1213/ANE.0000000000002741}.
Onnasch, L., Wickens, C.D., Li, H. and Manzey, D. (2013). Human
Performance Consequences of Stages and Levels of Automation. Human
Factors: The Journal of the Human Factors and Ergonomics Society, 56(3),
pp.476--488. doi: \url{https://doi.org/10.1177/0018720813501549}.
Sanchez, J.I. and De La Torre, P. (1996). A second look at the
relationship between rating and behavioral accuracy in performance
appraisal. Journal of Applied Psychology, 81(1), pp.3--10. doi:
\url{https://doi.org/10.1037/0021-9010.81.1.3}. Schmid, D. and Stanton,
N.A. (2019). Exploring Bayesian analyses of a small-sample-size
factorial design in human systems integration: the effects of pilot
incapacitation. Human-Intelligent Systems Integration, 1(2-4),
pp.71--88. doi: \url{https://doi.org/10.1007/s42454-020-00012-0}.
Schmutz, J.B., Outland, N., Kerstan, S., Georganta, E. and Anna-Sophie
Ulfert (2024). AI-teaming: Redefining collaboration in the digital era.
Current Opinion in Psychology, 58, pp.101837--101837. doi:
\url{https://doi.org/10.1016/j.copsyc.2024.101837}. Soetewey, A. (2020).
Correlation coefficient and correlation test in R. {[}online{]} Stats
and R. Available at:
\url{https://statsandr.com/blog/correlation-coefficient-and-correlation-test-in-r/}.
Suriya, S. and Nivetha, S. (2023). Interface Design - an overview
\textbar{} ScienceDirect Topics. {[}online{]} www.sciencedirect.com.
Available at:
\url{https://www.sciencedirect.com/topics/computer-science/interface-design}.
The R Foundation (2019). What is R? {[}online{]} R-project.org.
Available at: \url{https://www.r-project.org/about.html}. Vaccaro, M.,
Almaatouq, A. and Malone, T. (2024). When combinations of humans and AI
are useful: A systematic review and meta-analysis. Nature Human
Behaviour, {[}online{]} 8. doi:
\url{https://doi.org/10.1038/s41562-024-02024-1}. Wulff, J.N. and
Taylor, L. (2024). How and why alpha should depend on sample size: A
Bayesian-frequentist compromise for significance testing. Strategic
organization. doi: \url{https://doi.org/10.1177/14761270231214429}.
Yang, Z., He, L. and Huang, T. (2024). Powering Pro-Environment
Behavior: The Impact of Unlocking Reward Strategy on Pro-Environmental
Behavior. Sustainability, {[}online{]} 16(21), p.9561. doi:
\url{https://doi.org/10.3390/su16219561}.

\end{document}
